%  LaTeX support: latex@mdpi.com 
%  For support, please attach all files needed for compiling as well as the log file, and specify your operating system, LaTeX version, and LaTeX editor.

%=================================================================
\documentclass[jimaging,article,accept,pdftex,moreauthors]{Definitions/mdpi} 

%=================================================================
%----------
% journal
%----------

% MDPI internal commands
\firstpage{1} 
\makeatletter 
\setcounter{page}{\@firstpage} 
\makeatother
\pubvolume{1}
\issuenum{1}
\articlenumber{0}
\pubyear{2025}
\copyrightyear{2025}
\externaleditor{Firstname Lastname} % More than 1 editor, please add `` and '' before the last editor name
\datereceived{14 June 2025} 
\daterevised{16 July 2025} % Comment out if no revised date
\dateaccepted{21 July 2025} 
\datepublished{ } 
%\datecorrected{} % For corrected papers: "Corrected: XXX" date in the original paper.
%\dateretracted{} % For retracted papers: "Retracted: XXX" date in the original paper.
\hreflink{https://doi.org/} % If needed use \linebreak

%=================================================================
% Add packages and commands here. The following packages are loaded in our class file: fontenc, inputenc, calc, indentfirst, fancyhdr, graphicx, epstopdf, lastpage, ifthen, lineno, float, amsmath, setspace, enumitem, mathpazo, booktabs, titlesec, etoolbox, tabto, xcolor, soul, multirow, microtype, tikz, totcount, changepage, attrib, upgreek, cleveref, amsthm, hyphenat, natbib, hyperref, footmisc, url, geometry, newfloat, caption

\graphicspath{{figures/}} % path to folder with figures
%\usepackage{subcaption}
\usepackage{listings}
\usepackage{xcolor}
\usepackage{gensymb} % for \textdegree
\usepackage{tabularx}
\newcolumntype{Y}{>{\centering\arraybackslash}X}
\usepackage{amsmath,mathtools,amsthm}
	\setcounter{MaxMatrixCols}{15}
\usepackage{array}
\usepackage{amssymb}
\usepackage{amsfonts}
\usepackage{float}
\usepackage{chngcntr}
\counterwithin*{equation}{section}
\usepackage[super]{nth}
\usepackage{longtable}
\usepackage{tabularx}
\usepackage{multirow}
\usepackage{adjustbox}
\usepackage{rotating}
\usepackage{makecell}
	\renewcommand\cellalign{ll}

\definecolor{deepblue}{rgb}{0,0,0.5}
\definecolor{deepred}{rgb}{0.6,0,0}
\definecolor{deepmagenta}{RGB}{195,12,214}
\definecolor{deepgreen}{RGB}{29,122,30}
\definecolor{mintgreen}{RGB}{192,249,192}
\definecolor{codepurple}{RGB}{204, 0, 153}
\definecolor{LightCyan}{rgb}{0.88,1,1}
\definecolor{GainsboroPurple}{RGB}{247,176,247}
\definecolor{LightSlateGrey}{RGB}{124,118,118}
%    
\lstdefinestyle{mystyle}{
    commentstyle=\color{LightSlateGrey},
    keywordstyle=\color{deepgreen}\bfseries,
    emphstyle=\color{deepred},
    numberstyle=\tiny\color{deepblue},
    stringstyle=\color{codepurple},
    basicstyle=\ttfamily\scriptsize,
    breakatwhitespace=false,         
    breaklines=true,                 
    captionpos=b,                    
    keepspaces=true,                 
    numbers=left,                    
    numbersep=5pt,                  
    showspaces=false,                
    showstringspaces=false,
    showtabs=false,                  
    tabsize=2
}
\lstset{style=mystyle}

\usepackage{verbatim}

%=================================================================
% Full title of the paper (Capitalized)
\Title{Reclassification Scheme for Image Analysis in GRASS GIS Using Gradient Boosting Algorithm: A Case of Djibouti, East Africa}

% MDPI internal command: Title for citation in the left column
\TitleCitation{Reclassification Scheme for Image Analysis in GRASS GIS Using Gradient Boosting Algorithm: A Case of Djibouti, East Africa}

% Author Orchid ID: enter ID or remove command
\newcommand{\orcidauthorA}{0000-0002-5759-1089} % Polina Add \orcidA{} behind the author's name

% Authors, for the paper (add full first names)
\Author{Polina Lemenkova %MDPI: Please carefully check the accuracy of names and affiliations. % Author's reply: yes, correct.
 \orcidA{}}

%\longauthorlist{yes}

% MDPI internal command: Authors, for metadata in PDF
\AuthorNames{Polina Lemenkova}

% MDPI internal command: Authors, for citation in the left column
\AuthorCitation{Lemenkova, P.}
% If this is a Chicago style journal: Lastname, Firstname, Firstname Lastname, and Firstname Lastname.

% Affiliations / Addresses (Add [1] after \address if there is only one affiliation.)
\address[1]{%
  Department of Biological, Geological and Environmental Sciences, Alma Mater Studiorum %MDPI: We arranged the authors’ address information from subordinate to superior. Please check and confirm this modification. % Author's reply: yes, I confirm.
-%MDPI: We changed -- to -, please confirm % Author's reply: yes, I confirm.
Università di Bologna, Via Irnerio 42, Emilia-Romagna, 40126 Bologna, %MDPI: Please confirm if it can be removed. We can keep it if it is province abbreviation % Author's reply: not needed, I deleted "IT".
 Italy; polina.lemenkova2@unibo.it  Tel.: +39-3446928732 %MDPI: Please check if need add hyphen in it % Author's reply: yes, telephone mobile number is correct.
}

% Contact information of the corresponding author
%\corres{Correspondence:  (P.L.)}

% Current address and/or shared authorship
%\firstnote{Current address: Affiliation 3.} 

% Abstract (Do not insert blank lines, i.e. \\) 
\abstract{Image analysis is a valuable approach in a wide array of environmental applications. Mapping land cover categories depicted from satellite images enables the monitoring of landscape dynamics. Such a technique plays a key role for land management and predictive ecosystem modelling. Satellite-based mapping of environmental dynamics enables us to define factors that trigger these processes and are crucial for our understanding of Earth system processes. In this study, a reclassification scheme of image analysis was developed for mapping the adjusted categorisation of land cover types using multispectral remote sensing datasets and Geographic Resources Analysis Support System (GRASS) Geographic Information System (GIS) software. The data included four Landsat 8--9 %MDPI: Please check if it should be 8--9, please check it in whole text % Author's reply: yes, I corrected to "8--9" here and in the text.
 satellite images on 2015, 2019, 2021 and 2023. The sequence of time series was used to determine land cover dynamics. The classification scheme consisting of 17 initial land cover classes was employed by logical workflow to extract 10 key land cover types of the coastal areas of Bab-el-Mandeb Strait, southern Red Sea. Special attention is placed to identify changes in the land categories regarding the thermal saline lake, Lake Assal, with fluctuating salinity and water levels. The methodology included the use of machine learning (ML) image analysis GRASS GIS modules `r.reclass' for the reclassification of a raster map based on category values. Other modules included `r.random', `r.learn.train' and `r.learn.predict' for gradient boosting ML classifier and `i.cluster' and `i.maxlik' for clustering and maximum-likelihood discriminant analysis. To reveal changes in the land cover categories around the Lake of Assal, this study uses ML and reclassification methods for image analysis. Auxiliary modules included `i.group', `r.import' and other GRASS GIS scripting techniques applied to Landsat image processing and for the identification of land cover variables. The results of image processing demonstrated annual fluctuations in the landscapes around the saline lake and changes in semi-arid and desert land cover types over Djibouti. The increase in the extent of semi-desert areas and the decrease in natural vegetation proved the processes of desertification of the arid environment in Djibouti caused by climate effects. The developed land cover maps provided information for assessing spatial--temporal changes in Djibouti. The proposed ML-based methodology using GRASS GIS can be employed for integrating techniques of image analysis for land management in other arid regions of Africa.}

% Keywords
\keyword{\textls[-15]{image processing; image analysis; remote sensing; satellite image; geoinformatics;} machine learning}

% The fields PACS, MSC, and JEL may be left empty or commented out if not applicable
%\PACS{J0101}
%\MSC{}
%\JEL{}

\begin{document}

%----------- section ----------->
\section{Introduction}\label{sec1}

\subsection{Background}

{Satellite images play a pivotal role in environmental monitoring for reconstructing the} the history of landscapes on Earth. Understanding how land cover types change over time gives{ }insight into {the development of }human--environment interactions \cite{Ndehedehe2022,Foli,Speranza}. The main threat of environmental {changes are} linked to climate issues such as desertification, erosion, landscape fragmentation, coastal degradation {and} deforestation. Most of this information {can be discovered on spaceborne images due to the fundamental principle of {remote sensing (RS)}---{differences in spectral reflectance of pixels corresponding to diverse objects on the Earth.} %EE: please check intended meaning is retained.
Using such benefits of RS, many} environmental processes and phenomena can be recognised on satellite images {and Earth Observation (EO) data}, including, for instance, deforestation, soil erosion and desertification. 

The important application of {EO} data for environmental monitoring consists of time series analysis. Such a technique is based on the comparison of changes visible in several multi-temporal images to evaluate trends in vegetation coverage during a given temporal \mbox{period \cite{Springer2023,10281679}}. {This enables the detection of land cover changes and} the identification of {landscape} \mbox{dynamics \cite{su13084092,NATARAJAN2016256}}. EO data processing requires advanced techniques that allow us to find optimal and time-efficient solutions for solving complex geospatial tasks in image analysis. While traditional methods {of Geographical Information Systems (GISs) }have limitations, programming algorithms have significant benefits in the time series analysis of satellite images.

Technologies {of} {image analysis} play an essential role in the development {of environmental monitoring}. {Traditional tools for processing RS data rely on} {GISs} \cite{BORANA2023261,Zhu02012021}. However, {due to the large volume of RS data, the automation of image analysis is needed to effectively perform time series analysis and extract geoinformation}. {Recent {technological} %EE: please check intended meaning is retained.
developments have enabled novel methods for extracting geoinformation from spaceborne data, including} image enhancement \cite{Dyson2024}, image differentiation, rationing or principal component \mbox{analysis \cite{7877513}}. Machine learning (ML) methods have been found to improve the efficiency of {image analysis} through automated algorithms. Existing cases of ML algorithms in cartography and RS tend to be {related} to problems of image processing that are similar to those faced by decision algorithms{. }In each case, ML algorithms {aim to identify} {the best} solution strategy for image classification \cite{FERCHICHI2024122211}.

{Examples of} ML algorithms include{ }Random Forest (RF) \cite{Balha2022,Lemenkova202401}, {Support Vector Machines (SVMs)} \cite{10450641} and {Convolutional Neural Networks (CNNs)} \cite{SUAREZFERNANDEZ2024}. Such approaches are derived from programming and are adapted to{ }image analysis{ }\cite{HUANG2020104580,Arabameri}. Examples of {novel} techniques {in image analysis} include heuristic clustering \cite{8524429,10615674,669164,4666175}, flexible control of \mbox{variables \cite{8999627}}, optimised pattern recognition \cite{9708387,1634653}, ensemble learning \cite{10145147} and objective analysis \cite{8938484,9279780}. 

{Following such applications of ML in Earth and climate sciences, }this study {reports the case of a novel independent reclassification scheme for image analysis at the pixel level using} ML {coupled with} %EE: please check intended meaning is retained.
Geographic Resources Analysis Support System (GRASS) GIS and gradient boosting algorithm. {Its advantages include the optimisation approach, which} considers previous iterative cycles of image classification for the generation of {new} %EE: please check intended meaning is retained.
models and the narrowing of pixel assignment options to various land cover classes. 


%----------- section ----------->
\subsection{Objectives and Goals}

The main goal{ }is to{ identify}{ changes in land cover types }of Djibouti using the {time series of }satellite images {by the reclassification scheme}. {GRASS GIS version 8.4.1 %MDPI: Please state the version number of the software. All software used needs to be added with version number, please check it in the full text % Author's reply: yes, I inserted version of software here and in the text (Methodology section).
is an advanced scripting cartographic and RS software with} a powerful computational engine for geospatial processing \cite{NETELER2012124}. {Developed by the international GRASS developer community} \cite{grass_development_team_2025_14918754}, {it is used for processing cartographic data, maps, models and satellite image analysis across various applications in geoinformatics. Some of its important benefits are a built-in framework and a Python Application Programming Interface (API), which supports the automation of data processing through scripting.} To achieve this goal, the objective{ }is to map land cover changes in the coastal region of Djibouti{ }using RS data.{ }The{ }scripts are included and commented for{ }reference{. }The specific objectives of this research are as follows:

\begin{enumerate}[leftmargin=2.1em,labelsep=0.3mm]
	\item To produce land cover maps based on{ }the classification of Landsat satellite images {from} {2015}, 2019, 2021 and 2023; 
	\item To quantify and visualise dynamic changes in land cover types {for the identification of landscape dynamics}; 
	\item To assess the accuracy of{ }classification for vegetation mapping in{ }central and coastal zones {around} Lake Assal. 
\end{enumerate}

{The rest of this manuscript is organised as follows. Section \ref{sec1} presents the introduction with a related literature review and identifies the objectives and goals. {Section \ref{sec2} describes the study area, focusing on environmental and climate--hydrological settings that {influence} landscape types.} %EE: please check intended meaning is retained.
Section \ref{sec3} introduces the materials and numerical methods used in this study, {including the dataset and workflow, which are described as the methodological approach.} %EE: please check intended meaning is retained.
Major technical procedures are described in the subsections that follow. In Section \ref{sec4}, we present and discuss the results with cases of land cover changes, comment on uncertainties, and provide and interpret a series of maps. Finally, Section \ref{sec5} concludes the manuscript, summarising the findings and providing recommendations for future studies.}

%----------- section ----------->
\section{Study Area}\label{sec2}

This study area is located in Djibouti and surrounding coastal areas in the Gulf of Tadjoura, Bab-el-Mandeb Strait{ }(Figure \ref{fig01}). Situated in the region of the Horn of Africa and surrounded by waters of the southern Red Sea, Djibouti is a strategic point with geo-maritime importance{ }\cite{McCabe02102019,Lawale02102023}. Djibouti is one of the hottest and driest places on {Earth} with many issues related to climate and weather processes \cite{Verdugo2016}. {The} thermal crater saline lake, Lake Assal, {is located in this area, which is} the lowest point on land in East Africa with extreme levels of evaporation. {The topography affects }temperature and precipitation, {which }vary according to the elevation and differ in coastal and mountainous regions \cite{Haile}. 



This region experiences {fluctuations} in salinity and water levels {due to} seasonal and annual changes in temperature and precipitation. {Dominating landscapes are represented by }dry mountain vegetation {typical of arid climates}. The distribution and genera of flora and fauna are{ }controlled by low precipitation and high temperatures in semi-desert and desert environments \cite{GEBREMESKELHAILE2020135299,Lemenkova202221,AGHAKOUCHAK2015127}. For example, the 2011 drought in East Africa {affected rain-fed agriculture in} Djibouti, {which led to a regional }famine \cite{AGUTU2017287,DERCOLE2024104264,Dinku10112011}. Drought-induced crises endanger long-term investments and development,{ with} catastrophic human costs and effects on vulnerable regions of the country \cite{MISHRA2010202,FAVA202144,BRAVAR1991281}.

The drought-prone regions of{ }Djibouti {have experienced} strong effects from climate change, {including} fluctuations in the hydrological regime of thermal lakes. The arid climate{ has }triggered environmental problems{; for example,} it has desiccated and abandoned{ }agricultural lands, and salinated and blocked water sources due to sea level rise {and} flash floods \cite{Schulman02012019,CARRAO2016108}. Soil erosion and land degradation vary {across} topographically diverse landscapes, {reflecting the prevailing climatic conditions.} %EE: please check intended meaning is retained. 
\cite{FENTA2020135016,Knight2022}. Moreover, environmental pollution has been reported {along the coasts of the} Gulf of Tadjoura \cite{BOLDROCCHI2018812,BOLDROCCHI2018812,GAJDZIK2021112244}. 



{The} surroundings of the Horn of Africa are notable for climate--hydrological fluctuations which cause land cover changes \cite{Billi2022,SEKA2024141552}. The coastal region located in the Bab-el-Mandeb region includes diverse types of xeric grasslands and shrubland typical for arid ecoregions (Figure \ref{fig02}). {This region consists} of various land cover classes typical for semi-desert areas: dry shrub steppe and grassland, including species of \emph{Vachellia tortilis}, \emph{Vachellia nubica}, and \emph{Balanites aegyptiaca}. Many scarce vegetation spots include scattered trees or shrubs on the sandy plains of bare lands and semi-deserts.  

\begin{figure}[H]
   
	\includegraphics[width=12.0cm]{Fig_01.jpg}
  
    \caption{Topographic %MDPI: Please note that changes to the position/size of figures or tables may occur during the production stage. % Author's reply: yes, I confirm and agree.
 map of Djibouti with indicated study area (yellow rotated square) which shows the location of the Landsat images. Mapping software: Generic Mapping Tools (GMT) version 6.4.0. Data source: General Bathymetric Chart of the Oceans (GEBCO). Map source: Author.}
    \label{fig01}
\end{figure}

\vspace{-5pt}
\begin{figure}[H]
   
	\includegraphics[width=13.5cm]{Fig_02.jpg}
  
    \caption{\hl{Map} %MDPI: Please confirm whether the overlapping/incomplete content in this figure affects scientific understanding and if it does, please revise it. The contents of this figure are not legible. Please replace the image with one of a sufficiently high resolution (min. 1000 pixels width/height, or a resolution of 300 dpi or higher). Please change the hyphen (-) between numbers into an en dash (−, “U+2013”), e.g., “2019-2022” should be “2019–2022” and “0-9” should be “0–9”.
 of land cover types over Djibouti, showing the distribution of major vegetation and land cover types. Software: {Quantum Geographic Information System (QGIS)}. Data source: Food and Agriculture Organisation (FAO). Map source: Author.}
    \label{fig02}
\end{figure}

The eastern strip along the Red Sea coast is part of the Eritrean coastal desert \cite{TESFAI2021104327}. It includes vegetation along the muddy wadi mouths in diverse habitats \cite{MEDIN201526}, \mbox{grasslands \cite{ALI20231,IDOWU2020e00402}}, and secondary growth areas in the urban areas. The floodplains along the wadi rivers include scarce vegetation typical for semi-desert ecosystems. Xeric vegetation, farmland, mangroves along the coasts of the  Red Sea and valleys of wadi \cite{BLANCOSACRISTAN2022157098,BEYIN2024100247}, coral reefs along the coasts \cite{Klaus2015,COWBURN2019182,GAPPER2024e38249} and urban areas {collectively make up the region's landscape}. %EE: please check intended meaning is retained.
{ }Djibouti is one of the most{ seismically unstable }zones of Africa {with high} tectonic--magmatic activities \cite{Chandrasekharam,MOUSSA2023104765,AYELE2024105236,JALLUDIN1994237} due to the closeness of the Afar Triple Junction \cite{RIME2023104519,GIDAFIE2024102037,DARRAH201316}. The thermal regime around the Red Sea{ }is explained{ by its structural connection to the active{ }zone of the East African Rift} %EE: please check intended meaning is retained.
\cite{Varet2022,Lemenkova202206,PINZUTI2010169,Lemenkova202205,MESHESHA2021e08634}. 

{Such} complex geological setting and arid environment {have resulted in the} deficit of water {and favoured} the distribution of thermal springs with high salinity, %{of the brines} 
located in three regions---Lake Assal, Lake Hanle and Lake Abhe \cite{MOLOGNI2021107896,DAMORE1998197}. The craton origin of Lake{ Assal }is related to the Afar Triple Junction \cite{HOUSSEIN2010220,AWALEH2024102894,Houssein02112014}. The {geochemical composition of the} Assal Rift {includes }evaporitic halite deposits \cite{CASTANIER1992811,GASSE198967}. Its location in the arid zone with strong evaporation and no riverine input contributed to the  lake's increase in salinity, with a high brine concentration above 100 g/L \cite{DEKOV2021120126}. {Fluctuations of this concentration} %EE: please check intended meaning is retained.
are related to evaporation, the hot climate and geothermal processes \cite{ZAN1990561}. 

%----------- section ----------->

\section{Materials and Methods %MDPI: In Section Materials and Methods, if there is manufacturer, please state the name of the manufacturer, city, and country from where the equipment was sourced. % Author's reply: checked and corrected for software (added version, producer, city, and country).
}\label{sec3}
Our methodology presents a case of image reclassification and analysis using {ML} techniques{ }for environmental mapping in Djibouti, East Africa. A classification scheme consisting of 17 initial land cover classes was employed through a logical sequence of GRASS GIS modules to extract land cover information of the study area in Djibouti. {Reclassification includes the sequence of changing the values in a raster dataset to new values. The process is based on a set of rules which defined the criteria of updates in classes. This operation in GRASS GIS was applied for the simplification of land cover polygons, standardisation of assignment of pixels to classes and preparation of data for land cover analysis.}

%----------- subsection ----------->
\subsection{Data}

The data used in this study include {four images from the }Landsat 8--9  Operational Land Imager (OLI) and Thermal Infrared Sensor (TIRS): {the training image from 2015, {and images with biannual intervals from }2019, 2021 and 2023 (Figure \ref{fig03}). }%EE: please check intended meaning is retained.



The {Landsat} images were selected as the data source due to {their} quality and {reliability}\cite{GONG2024149,ALLU2025127478,CHANUWANWIJESINGHE2024112681}. {Earlier studies reported the effective} use of {the} RS data {in environmental studies which prove their utility and value }\cite{WANG2023311,DAREM2023341,FAROOQ2024101168,AFUYE2024103262,CHUGHTAI2021100482,VELASTEGUIMONTOYA2023100659,FARHAN2024103689}. {Nowadays}, satellite images are {excellent} sources of {geoinformation} {for mapping} arid regions \cite{AMAZIRH2023104586,FERCHICHI2024122211,MEKONNEN20251}. {Spaceborne data are successfully} used {in} the analysis of land cover change and deforestation in tropical forests, and{ for reporting }landscape {dynamics} \cite{PARRAPAITAN2024103796,ZHAO2024506,NZABARINDA2025114849,MONTEROBOTEY2024e03068}. {The} importance of Landsat {data} for environmental {research} has been described earlier \cite{AKSOY20221072,NESARI2024101989,KUMAR2022100578}.

In this study, four Landsat images were taken on the following dates: (1) 02 July 2015; (2) 23 July 2017; (3) 19 August 2021; and (4) 16 July 2023. The metadata{ are }summarised in Table \ref{tab01}. The cloudiness percentage is less than 10\% for all the images. The{ }technical parameters of {the} scenes are provided in Table \ref{tab02}. Common data for all four images are as follows: Datum and Ellipsoid: {{World Geodetic System 1984 (WGS84)}; Product Map Projection L1: Universal Transverse Mercator (UTM) zone 38 (for Djibouti);  Worldwide Reference System (WRS): Path--166, Row--52; Landsat Collection Number: 2; Category: T1; Station Identifier: LGN; and Sensor Identifier: OLI/TIRS. }%EE: please check intended meaning is retained.
The images were taken during day time in nadir with a roll angle of 1\degree  for the 2015 image and 0\degree for the images from 2017, 2021 and 2023. 

\begin{figure}[H]
  
	\hspace{+2pt}\includegraphics[width=13.0cm]{Fig_03.jpg}
  
    \caption{Original Landsat 8--9 OLI/TIRS scenes of Djibouti from 2015, 2019, 2021 and 2023 used for image processing. Image source: United States Geological Survey (USGS). Compilation: Author.}
    \label{fig03}
\end{figure}


\begin{table}[H]\small
\caption{Identifiers (IDs) and acquisition dates of the multispectral Landsat 8--9 OLI/TIRS images. %MDPI: It is same as table 3 caption, please check if need keep % Author's reply: the table caption is corrected.
}
\label{tab01}
\setlength\tabcolsep{3pt}
\begin{tabularx}{\textwidth}{ X l X }
\toprule
\textbf{Date} & \textbf{\hspace{-26pt}Landsat Product Identifier} & \textbf{Scene ID} \\
\midrule
	02 July 2015 & \hspace{-26pt}LC08\_L2SP\_166052\_20150702\_20200909\_02\_T1 & LC81660522015183LGN01 \\
	23 July 2017 &  \hspace{-26pt}LC08\_L2SP\_166052\_20170723\_20200903\_02\_T1 & LC81660522017204LGN00 \\
	19 August 2021 & \hspace{-26pt}LC08\_L2SP\_166052\_20210819\_20210827\_02\_T1 & LC81660522021231LGN00 \\
	16 July 2023 & \hspace{-26pt}LC09\_L2SP\_166052\_20230716\_20230719\_02\_T1 & LC91660522023197LGN00 \\
\bottomrule
\end{tabularx}
\end{table} 

The{ image quality score was 9 }%EE: please check intended meaning is retained.
for all scenes, according to the Landsat specifications. The remaining various metadata for each scene are summarised in Table \ref{tab02}. The geographic characteristics of the images are referenced in Table \ref{tab03}. {The Landsat Scene Identifiers (IDs) are \mbox{provided below.} %EE: please check intended meaning is retained.
}



 The images have been filtered to reduce noise due to cloudiness and atmospheric effects. Images were systematically selected from the USGS EarthExplorer repository and taken at the end of the dry season. During this period, cloud cover is generally low but vegetation appears bright green and contrasts with all other land cover types. 


\begin{table}[H]
\caption{{Metadata%MDPI: We removed the vertical lines in all tables. Please confirm. % Author's reply: yes, I confirm.
} of the Landsat 8--9 OLI/TIRS satellite images{. }Source: USGS (EarthExplorer).}
\label{tab02}
    \centering
\footnotesize 
    \begin{tabular}{ l  l  l  l  l }
 \toprule
        \textbf{Dataset Attribute} & \textbf{Attribute Value} & \textbf{Attribute Value} & \textbf{Attribute Value} & \textbf{Attribute Value} \\ \midrule
        Date Acquired & 16 July 2023 %MDPI: Please revise date format to  Date Month Year, e.g.,  1 January 2020. The same for the following highlights. % Author's reply: corrected.
 & 19 August 2021 & 23 July 2017 & 02 July 2015 \\ \midrule
        Date Product Generated L2 & 19 July 2023 & 27 August 2021 & 03 September 2020 & 09 September 2020 \\ \midrule
        Date Product Generated L1 & 16 July 2023 & 27 August 2021 & 03 September 2020 & 09 September 2020 \\ \midrule
        Land Cloud Cover & 0.02 & 0.19 & 0.45 & 0.08 \\ \midrule
        Scene Cloud Cover L1 & 0.06 & 0.71 & 12.03 & 0.08 \\ \midrule
        Ground Control Points Model & 752 & 722 & 742 & 768 \\ \midrule
        Ground Control Points Version & 5 & 5 & 5 & 5 \\ \midrule
        Geometric RMSE Model & 3563 & 2963 & 2977 & 2098 \\ \midrule
        Geometric RMSE Model X & 2240 & 2009 & 1967 & 1089 \\ \midrule
        Geometric RMSE Model Y & 2771 & 2179 & 2235 & 1794 \\ \midrule
        Processing Software Version & LPGS\_16.3.0 & LPGS\_15.5.0 & LPGS\_15.3.1c & LPGS\_15.3.1c \\ \midrule
        Sun Elevation L0RA & 62.48488638 & 64.48715625 & 62.85521856 & 62.40770766 \\ \midrule
        Sun Azimuth L0RA & 65.62042782 & 84.79256539 & 68.51197578 & 61.89003236 \\ \midrule
        TIRS SSM Model & N/A & FINAL & FINAL & ACTUAL \\ \midrule
        Scene Center Latitude & 11.56786 & 11.56779 & 11.56783 & 11.56778 \\ \midrule
        Scene Center Longitude & 42.94139 & 42.97394 & 42.97879 & 42.96942 \\ \midrule
        Corner Upper Left Latitude & 12.60889 & 12.60924 & 12.60931 & 12.60921 \\ \midrule
        Corner Upper Left Longitude & 41.88103 & 41.91137 & 41.91688 & 41.90861 \\ \midrule
        Corner Upper Right Latitude & 12.62518 & 12.62530 & 12.62532 & 12.62529 \\ \midrule
        Corner Upper Right Longitude & 43.98441 & 44.01755 & 44.02308 & 44.01479 \\ \midrule
        Corner Lower Left Latitude & 10.50122 & 10.50151 & 10.50157 & 10.50149 \\ \midrule
        Corner Lower Left Longitude & 41.90427 & 41.93439 & 41.93986 & 41.93165 \\ \midrule
        Corner Lower Right Latitude & 10.51472 & 10.51483 & 10.51484 & 10.51482 \\ \midrule
        Corner Lower Right Longitude & 43.99199 & 44.02489 & 44.03037 & 44.02215 \\\bottomrule
    \end{tabular}
\end{table}

\vspace{-6pt}
\begin{table}[H]
\caption{Technical characteristics and geospatial location of the multispectral Landsat images.}\label{tab03}
\setlength\tabcolsep{0pt}
\begin{tabularx}{\textwidth}{ X X X X}
\toprule
\textbf{N-S} & \textbf{Data Values} & \textbf{E-W} & \textbf{Data Values}\\
\midrule
	North: & 1,%MDPI: We added commas to separate thousands for numbers with five or more digits. Please confirm. % Author's reply: yes, I confirm.
395,915.00 &  East: & 393,015.00 \\
	South: &  1,162,485.00 &  West: & 164,085.00 \\
	Resolution: &  30.00 m & Cells: & 59,376,811 \\
	Rows: & 7781 & Columns: & 7631 \\
\bottomrule
\end{tabularx}
\end{table}

%----------- subsection ----------->
\subsection{Workflow}

{The schematic workflow of the research is presented in Figure} \ref{fig04}.
 


The{ }image processing used to create a land cover map from Landsat OLI/TIRS imagery was carried out in six steps using {the following techniques of} GRASS GIS: 
\begin{itemize}
	\item Image import;
	\item Data exploration and analysis;
	\item Classification of images using MaxLike method;
	\item Reclassification scheme using `echo' of Unix and `r.reclass' module;
	\item Classification of images using gradient boosting algorithm of ML;
	\item Image post-processing and mapping;
	\item Accuracy assessment. 
\end{itemize}

\begin{figure}[H]
  
	\hspace{+2pt}\includegraphics[width=13.0cm]{Fig_04.jpg}
  
    \caption{Workflow scheme. Source: Author.}
    \label{fig04}
\end{figure}

The {workflow is based on scripting techniques of} GRASS GIS {for image }analysis (software developed originally by the U.S. Army Construction Engineering Research Laboratories in 1982, now maintained by he developer team in Champaign, Illinois, U.S.A.). {The} topographic map is plotted using  existing techniques in GMT version 6.4.0.~\cite{WRIGHT1998737,Lemenkova202217,9921513,Lemenkova202203,10537274} (software developed originally in 1988 by Paul Wessel and Walter H. F. Smith, in Lamont-Doherty Earth Observatory of Columbia University, Palisades, NY, U.S.A.). 

%----------- subsection ----------->
\subsection{Image Preprocessing}

The Landsat 8--9 OLI/TIRS images were geometrically rectified according to EPSG coordinates: 2095, Universal Transverse Mercator (UTM) zone for Djibouti according to the USGS technical characteristics. In this study, a scripting method of GRASS GIS was applied, which uses several modules with code snippets explained and commented below. 

{The} images were imported into the working folder using the `r.import' module,{ then }the same procedure was repeated for all the bands and the four images used in this study (2015, 2019, 2021 and 2023). Afterwards, the relief map was imported from GEBCO{. }The region was set up to the geospatial extent {with the} isolines{ }obtained{ }using the `r.contour' module{ }with a 200 m interval{.} To facilitate data processing, the best combinations of Landsat bands were{ }stacked into {coloured composites}. The bands that produced the best separation for pixels {were processed, excluding} the panchromatic channels{. The} {Tag Image File Format (TIFF)} files were joined using the `i.group' module {and enhanced using preprocessing steps which included} contrast stretching{ and }atmospheric correction. 

%----------- subsection ----------->
\subsection{Clustering}

The clustering was performed by the `i.cluster' module of GRASS GIS, which generates a signature file and reports results using the k-means clustering algorithm, as explained in previous studies \cite{5930121,10430357,Lemenkova202315,5739741}.{ }{Signature file determines the categories (classes) of land cover types{ using the analysis of spectral reflectance of the pixels} and contains{ }clusters and covariance matrices for each image.} %EE: Please check intended meaning is retained.­
{The initial means for each class before data processing reported in Appendix} \ref{AppA} {indicate {the starting seed values} %EE: Please check intended meaning is retained.­
of the multispectral band before iteration is assigned according to the signature files. The means are then calculated  during the `i.cluster' procedure and recalculated following the principles of maximal class separation (the contrast between the classes) and minimal class size (the details of classification with respect to a spatial resolution of 30 m). The table in Appendix} \ref{AppA} {reports such values for each of the 10 classes and 7 multispectral bands of Landsat, starting from 1---`Coastal aerosol' to 7---`SWIR-2'.} 

{The} initial {mean} %EE: Please check intended meaning is retained.­
values for each land cover class for the years 2015, 2017, 2021 and 2023 are reported in Tables \ref{tabA1}--\ref{tabA4}; see Appendix \ref{AppA}. Afterwards, the assignment of pixels was performed {iteratively}  %EE: Please check intended meaning is retained.­
for each land cover class. The same was applied for all images for the years 2015, 2017, 2021 and 2023 with a biannual gap. The classification was performed using the maximum-likelihood discriminant analysis classifier by the `i.maxlik' module \cite{CHOWDHURY2024100800,Lemenkova202323}. {The algorithm used} the signature file{ }calculated in the previous step. {Using} the signature file, the land cover types{ were classified }using maximum-likelihood discriminant analysis (Figure \ref{fig05}).

 \begin{figure}[H]
 
	\includegraphics[width=13.50cm]{Fig_05.jpg}
  
    \caption{Maps %MDPI: 1. Please change the hyphen (-) into a minus sign ($-$, "U+2212"). e.g., "-1" should be "$-$1". 2. Please confirm whether the overlapping/incomplete content in this figure affects scientific understanding and if it does, please revise it. % Author's reply: I revised hyphen in the figure into minus sign. I confirm that the content in this figure is correct and does not affects scientific understanding.
 of land cover types based on the classified Landsat 8--9 OLI/TIRS images from 2015, 2019, 2021 and 2023 covering Djibouti. The approach involved the maximum-likelihood discriminant analysis classifier. Background relief: GEBCO. Software: GRASS GIS. Mapping source: Author.}
    \label{fig05}
\end{figure}

%----------- subsection ----------->
\subsection{Accuracy Analysis}
 
The analysis of accuracy (Figure \ref{fig06}) revealed the areas of misclassified pixels {within water bodies}, %EE: Please check intended meaning is retained.­ % Author's reply: yes, intended meaning is retained.
including the Bab-el-Mandeb Strait. The accuracy assessment, showing the probability of correctly assigning each pixel to its respective class across all images, is presented in Figure \ref{fig06} and calculated with the averages of each calculated band. Afterwards, based on the classified ten categories in the images, two of them {(`water in the sea' and `water in shelf' areas)} %EE: Please check intended meaning is retained.­ % Author's reply: yes, intended meaning is retained.
were merged as a single class `water', and the names of the classes were assigned to each numerical cluster. 

\vspace{+5pt}
 \begin{figure}[H]
  
	
\begin{adjustwidth}{-\extralength}{0cm}
\centering
\includegraphics[width=15.0cm]{Fig_06.jpg}
\end{adjustwidth}
  
    \caption{Accuracy %MDPI: We revised the hyphen in the figure into minus sign. Please confirm. Please confirm whether the overlapping/incomplete content in this figure affects scientific understanding and if it does, please revise it. % Author's reply: yes, I confirm the revised hyphen in the figure into minus sign. I confirm that the content in this figure is correct and does not affects scientific understanding.
 analysis of image classification using reject threshold probability algorithm and the chi-squared test, which was applied to the reclassified images. {Source}: Author.}
    \label{fig06}
\end{figure}

After this step, ML was applied to all the images. The ML step started with the generation of training pixels from land cover classification performed in 2015, which were used as seed data. The extraction of training data was conducted using the `r.random' module{. }Afterwards, {the} training pixels were used for classification on the Landsat images (example for 2023) by training a model using the `r.learn.train' module. {The} gradient boosting classifier was used as an algorithm to perform ML classification. {Next, the} prediction was performed using the ‘r.learn.predict' ML module,{ }which applied a fitted estimator from the Python's Scikit-Learn library \cite{pedregosa2018} to raster images in {the} imagery group.{ In the next step, }the computed raster categories were automatically applied to the classification output and checked using the ‘r.category' module{.}

%----------- subsection ----------->
\subsection{Gradient Boosting}

{The} images were{ }processed using the {gradient boosting ML {method of image} classification. The essential concept of this algorithm is setting the target outcomes of classified images for the next model based on previous iterations, which in the end minimises the error of classification. Gradient boosting presents an advanced ensemble technique {of ML, }which combines the predictions of multiple classification steps by the principle of decision trees. Such an approach presents a novel method of satellite image processing{ }based on the supervisor approach. {The} algorithm{ }classifies{ }images using the iterative approach which considers the previous and next classification steps as loops. {It} combines several decision trees as{ }simplified learners to create a strong prediction mode{. }By minimising a loss function embedded in the algorithm, each new learner{ fixes} the mistakes committed by the group of prior learners. In this way, the accuracy and predictive capacity{ of the classifier} are increased {iteratively}, which is directed by the gradient of the loss function. 

{The core idea is that} the algorithm finds the best possible next model when combined with previous {ones}. In this way, it minimises the overall prediction error. The advantage of such an approach is that it {improves the accuracy of classification} through the convergence of decision{ }cycles and assigns cells that constitute the mosaic of pixels on{ }a raster image to the defined {categories}. When several {images are processed as a time series}, this enables us to effectively monitor landscape dynamics and detect{ changes}. 

The class separability matrices were computed for each pair of years (2015--2017 and 2021--2023) and show the distinguishability between classes (Figures \ref{figSep1} and \ref{figSep2}). The iterative cycles of image classification are summarised in  Appendix \ref{AppB} in four statistical tables for each of the study periods, respectively. {They report the pixels' distribution by categories of land cover classes for each year: 2015, 2017, 2021 and 2023. The values indicate percent convergence,{ which shows at which values the cluster means of land cover categories become stable during iterative classification}. %EE: Please check intended meaning is retained.­
Hence, iterative processes aim to achieve the maximal percentage of pixels that no longer move from cluster to cluster during iteration and reach the best possible values.} 

%----------- subsection ----------->
\subsection{Reclassification}

{To improve} accuracy, the images were reclassified using the principle of reclassification scheme. {When clusters are being generated by the algorithm, the means of classified categories constantly change, because cells are assigned to these classes and the mean is then recalculated to include new pixels. The goal of the iteration is to maximise the distances between the classes through recalculation. To this end, after the creation of all clusters and the assignment of pixels into these groups, the algorithm ‘i.cluster' changes cluster means through iteration processes. In this way, the algorithm attempts to increase the contrast between the classes, i.e., the numerical gap in values between the classes.}

%---------------------------->
\subsection{{Generalisation of FAO Land Cover Map}}

{The {FAO} land cover map of Djibouti was generalised from the classes of the Land Cover Classification System (LCCS) according to a technical report on the country. {The scheme covering land {cover} categories {comprising} 10 classes was used and adopted for the study areas within the country} %EE: Please check intended meaning is retained.­
(Table} \ref{tabFAO}). This part of the work was performed using QGIS software, version \mbox{3.42.2 ‘M\"unster'}. The QGIS was initially developed by G. Sherman and now a project of the Open Source Geospatial Foundation. 

{The reclassified }image was generated for each year using the ‘r.reclass’ module of GRASS GIS {and} ‘echo’ Unix command.{ }The raster map was reclassified based on the category values, and new raster maps for the years 2015, 2019, 2021 and 2023 were {generated}. The category values of the reclassified maps are based on an {iteration} of the categories using the information generated in ‘landusereclass.txt' file. {These} categories were {controlled} using the ‘r.category' module. The maps were{ }visualised using{ }modules {`d.mon' and ‘g.region'}.  The color palette was {selected from the} color tables. {Afterwards, the maps} were displayed using a combination of ‘d.rast’ and ‘d.vect’ {modules}. {Finally, }a cartographic grid was added {along with} the legends by the ‘d.legend’ module.


\begin{table}[H]
\caption{{Training pixels for validation {under} {LCCS} {for} %EE: Please check intended meaning is retained.­
Djibouti, derived from FAO classification scheme of land cover types: spatial distribution of the training ground truth samples.}}\label{tabFAO}%
\small
\begin{tabularx}{\textwidth}{llll}
\toprule
\textbf{ID} & \textbf{Landsat 8--9 OLI/TIRS Land Cover Class} & \textbf{FAO LCCS Codes}  \\
\midrule

       1 & Salt Plains and Hardpans & 6,020 \\ \midrule
        2 & Barren Land & 6,001, 6,004  \\ \midrule
        3 & Water Bodies & 7,001 // 8,001  \\ \midrule
        4 & \makecell{Mangroves and Aquatic Vegetation on Regularly Flooded \\(Temporarily or Permanently) Fresh or Brakish Water} & \makecell{41,653%MDPI: Please check if need use commas to separate thousands for numbers with five or more digits (not four digits) in this table, e.g., "10000" should be "10,000". Please check and revise all if needed % Author's reply:  corrected to separated thousands.
-R1,R2 // 41,739-R1 \\ 41,521-R1,R2 // 41,597-R1,R2}  \\ \midrule
        5 & Bushes and Shrubs & 11,490 // 11,494 \\ \midrule
        6 & \makecell{Farmland, {Shrubs} %EE: Please check intended meaning is retained.­
        	and \\Irrigated Trees} & \makecell{11,499 // 11,500 // 30,001, \\11,491, 11,495}  \\ \midrule
        7 & Built-up Areas, Artificial Surfaces and Associated Areas & 0,010  \\ \midrule
        8 & Broadleaved Semi-deciduous Forest // Woodland & \makecell{21,496 // 21,497--15,048 // \\21,496--12,134 // 21,497-12,940} \\  \midrule
        9 & Mosaic Cropland (Cultivated Areas) & 0,003--0,004 \\ \midrule
        10 & Sparse Vegetation & 20,049, 20,056 20,058  \\
\bottomrule
\end{tabularx}
\noindent{\footnotesize{Information source: FAO (land use statistics).}}
\end{table}

The reclassification was performed since water surfaces in the Bab-el-Mandeb Strait and southern Red Sea had different colours. Nevertheless, they belong to the same class of water. The differences in the coastal regions were smooth at the time of image acquisition and appear as dark colours in all images. {With each iteration, the values of the {means shift to a higher} %EE: Please check intended meaning is retained.­
percentage of pixels within each cluster. This is illustrated in the class separability matrix, which shows the distinction between the categories} (\mbox{Figures \ref{figSep1} {and } \ref{figSep2}}; Appendix \ref{AppC}). {As means never become completely static, a \% convergence and a maximum number of iterations are defined to finalise the process of reclassification before the maximum number of cycles {is reached}. %EE: Please check intended meaning is retained.­
The final values of the computed class means for each land category are reported in Appendix} \ref{AppD}. {Once the maximum number of reclassification cycles is reached, the optimal \% convergence is achieved through the increased number of iterations during the ‘i.cluster' procedure. Here, the number of cycles depends on the complexity of the terrain and land cover patterns.} 

%\pagebreak
%----------- section ----------->
\section{Results and Discussion}\label{sec4}

\subsection{Land Cover Change Analysis}
{The results of the classified }satellite images are {presented in Figure} \ref{fig05}. The{ }study {area} was categorised into 10 distinct land cover classes following the FAO classification scheme: \mbox{(1) salt} plains and hardpans; (2) barren land; (3) water bodies; (4) mangroves and aquatic vegetation on regularly flooded (temporarily or permanently) fresh or brackish water; (5) bushes and shrubs; (6) farmland, {shrubs} %EE: Please check intended meaning is retained.­
and irrigated trees; (7) built-up areas, artificial surfaces and associated areas; (8) broadleaved semi-deciduous forest and woodland; (9) mosaic cropland in cultivated areas; and (10) sparse vegetation in desert areas. Among these classes, the following categories included sub-clusters: water bodies comprise ponds, areas covered by marshes, rivers, estuaries, and coastal waterways; bushland and shrubland include non-classified types of vegetation such as grasslands, shrubs, rangelands, and savannah; and coastal areas surrounding Lake Assal are categorised by salt-tolerant vegetation and mangroves.}



{{Over an estimated period of a decade, from 2013 to 2023,} %EE: Please check intended meaning is retained.­
image analysis revealed certain changes in land cover types across the study area of Djibouti. Thus, it is noteworthy that the amount of mangroves along the coastal areas decreased slightly, shrinking by 0.35 km\textsuperscript{2} from 6.28 km\textsuperscript{2} to 5.93 km\textsuperscript{2} during this period. Furthermore, a large area of 1392 km\textsuperscript{2} that had previously been covered by bushes and shrubland was converted into other land use categories, resulting in a decrease in coverage from 1194 km\textsuperscript{2}. As a result of this reduction, the percentage of bush cover decreased {to} %EE: Please check intended meaning is retained.­
15\%. Over the course of the studied period, {both salt-covered areas around Lake Assal and croplands} %EE: Please check intended meaning is retained.­
showed downward tendencies, with farmland areas falling from 3.21 km\textsuperscript{2} to 2.75 km\textsuperscript{2} and salt coverage from 17 km\textsuperscript{2} to 16.04 km\textsuperscript{2}.}

Additionally, saline waters mixed with the fresh waters of the estuaries of the Gulf of Tadjoura have a lighter coloration than those of the open sea. Vegetation types such as mosaic croplands, sparse vegetation, xeric shrubland and grassland or bare soil areas are detected and indicated on the maps. Broadleaved deciduous vegetation is relatively scarce in the study area. The regions of shrub, and flooded brackish water areas are mostly occupied by mangrove forests, which have stronger spectral signals and appear as various shades of colors. The computational results of the pixels by land cover classes are summarised in Table \ref{tab04}.

\begin{table}[H]
\caption{Distribution of pixels by land cover classes: biannual variations in study area of Djibouti over 10 land cover classes in summer periods from 2015 to 2023.}\label{tab04}

\begin{tabularx}{\textwidth}{ p{1.5cm}  X  X  X  X  X  X  X  X  X  X  }
  \toprule
   \textbf{Class} & \textbf{1} & \textbf{2} & \textbf{3} & \textbf{4} & \textbf{5} & \textbf{6} & \textbf{7} & \textbf{8} & \textbf{9} & \textbf{10} \\
   \midrule
    2015 & 1531 & 251 & 314 & 755 & 964 & 761 & 699 & 616 & 448 & 591 \\
    2017 & 1452 & 361 & 120 & 746 & 1031 & 813 & 672 & 714 & 502 & 609 \\
    2021 & 588 & 1225 & 223 & 699 & 1007 & 798 & 608 & 644 & 495 & 644 \\
    2023 & 842 & 915 & 367 & 733 & 965 & 840 & 672 & 627 & 533 & 547 \\
    \bottomrule
\end{tabularx}

\end{table}

Xeric shrubland and sparse deciduous forests typical in Djibouti {are} located along the coasts of the Red Sea or Gulf of Tadjoura. {Spectral} signatures {of this class} differ from those of mosaic vegetation (semi-deciduous forests or sparse lands) found in inland regions of the country, on the border with Ethiopia and Eritrea. {Thin} coastal regions covered with mangrove plants along the coasts of the Bab-el-Mandeb Strait {are identified as} wetlands partially submerged in water. This {type} of vegetation shows {little} difference compared to the salt areas of Lake Assal, which experience significant changes. 

{There were also notable expansions in other land cover categories of landscapes in Djibouti.  For example, the area of bare lands and in the interior deserts (western region of the country) increased significantly between the estimated period of 2015 and 2023. Specifically, the areas expanded from {5,56} %EE: Attention AE: please check formatting.
km\textsuperscript{2} to 6.24 km\textsuperscript{2}. This change highlights the processes of desertification of regional landscapes, which is caused by recent climate change and global warming. Hence, desertification processes have become notable over time, with arid regions occupying more areas in inner regions of the country.}

\subsection{Uncertainties and sources of error}

{Understanding possible uncertainties underlying the classification schemes in RS data processing is imperative. }Accurate vegetation mapping using satellite image classification is not easy to achieve due to the spectral confusion between different vegetation types and{ }similarity in spectral reflectance of various plant species \cite{1370179,Dobrowski,Moore,8688643}. {Therefore, accuracy assessment is a crucial step in the evaluation process. In this case, the generated map was compared to the ground truth values, represented in this case by sample points taken from the FAO-based map}. {The overall accuracy (OA) and Kappa coefficient {(KC)} were computed to evaluate the classified map's accuracy by measuring the similarity between the sample points and the correctly and incorrectly categorised classes (Table} \ref{tabKappa}). 

\begin{table}[H]
\caption{{Error matrix for validation of gradient boosting ML algorithm {across} land cover categories.}}
\label{tabKappa}

\begin{adjustwidth}{-\extralength}{0cm}
\begin{tabularx}{\fulllength}{ p{0.7cm}  p{0.9cm}  p{0.9cm}  p{0.9cm}  p{1.5cm}  p{0.9cm}  p{1.3cm}  p{1.3cm}  p{1.3cm}  p{1.5cm}  p{1.5cm}  p{0.7cm} }
  \toprule
  Class & \makecell{Salt \\Plains} %MDPI: Please check whether the font direction of the table header should be maintained  % Author's reply: I corrected the font direction of the table into normal position (no rotation).
  & \makecell{Barren \\ Land} & \makecell{Water} & Mangrove & Shrub & Farmland & \makecell{Built-Up \\ Area} & Forest & Cropland & \makecell{Sparse \\ Vegetation} & Total \\
   \midrule
    {1} & {3503} & 12 & 4 & 7 & 31 & 4 & 29 & 2 & 9 & 21 & 3622\\  
    {2} & 1 & {1651} & 43 & 4 & 5 & 0 & 8 & 12 & 0 & 2 & 1726  \\
    {3} & 0 & 1 & {1018} & 21 & 0 & 8 & 15 & 4 & 2 & 12 & 1081  \\  
    {4} & 2 & 3 & 14 & {945} & 1 & 0 & 38 & 51 & 27 & 12 & 1093  \\ 
    {5} & 3 & 18 & 256 & 12 & {2751} & 12 & 48 & 35 & 0 & 0 & 3135 \\
    {6} & 22 & 16 & 39 & 0 & 2 & {546} & 5 & 8 & 31 & 83 & 752  \\
     {7} & 2 & 6 & 10 & 1 & 12 & 14 & {1132} & 44 & 6 & 13 & 1240  \\
     {8} & 1 & 1 & 0 & 204 & 37 & 52 & 0 & {1439} & 2 & 0 & 1736 \\
     {9} & 3 & 12 & 2 & 62 & 14 & 6 & 0 & 45 & {1968} & 0 & 2033 \\
     {10} & 14 & 2 & 0 & 16 & 83 & 12 & 0 & 26 & 0 & {2382} & 2535 \\ 
     \midrule 
     {} & \multicolumn{11}{c}{{OA} = 93.4\%; {KC} = 87.5\%} \\ 
    \bottomrule
\end{tabularx}
\end{adjustwidth}
\end{table}



{The sources of errors in automatic classification are also caused by the individual patterns of land cover categories on Earth; the }data are not normally distributed in most cases. {Moreover, }the light signature of the deciduous broadleaved plants enabled us to detect occasionally growing plants on the mountainous slopes in the central regions of the country where precipitation is higher than in semi-desert areas. {Thus}, deciduous shrubland in all the images is attributed to the high backscatter coefficients of this vegetation type and spectral reflection effects. Therefore, validation using the computed {KC} was performed and is reported in Table \ref{tabKappa}.

{In order to avoid uncertainties in classification caused by such cases,} coastal types of vegetation {were} grouped in the same class with other plants having similar spectral properties (grassland). {The} classification accuracy {was performed in order to evaluate possible sources of errors. To this end, images} were processed using a rejection probability test, which examines the correctness of the pixel’s classification using {the} chi-squared test. This made it possible to evaluate the classification of different types of land {cover types over the} study area using a traditional{ }approach of GRASS GIS based on automated clustering (Figure \ref{fig06}). 

\subsection{Interpretation and Data Analysis}

Major land cover types in Djibouti {include} xeric shrubland, bare land areas in mountainous deserts, estuary of the Gulf of Tadjoura, cropland areas and sparse vegetation{. Multispectral} bands of the images were analysed and compared using a spectral diagram and spatial analysis. The main land cover maps of Djibouti generated using{ reclassified scenes} and ML-based {analysis} for the images{ is displayed in }Figures \ref{fig07} and \ref{fig08}.



The reclassified images are shown in Figure \ref{fig07}. Here, the areas of water are merged into one class, which improved the classification. Nevertheless, the neighbouring regions containing vegetation types with similar pixel reflectance {were misclassified, which required} the ML approach for image processing. A fluctuation in Lake Assal visible on the images is related to climate effects and geothermal activities, as is also mentioned in previous studies. This study revealed that similar to the fluctuations of saline sabkhas in the northern Sahara \cite{Lemenkova202322,Lemenkova202313}, the saline lake, Lake Assal, also witnessed slight changes in the water surface and extent of the lacustrine coasts. Nevertheless, the nature of such physiographic variations is more related to the geologic origin \cite{Eldosouky2021}. 

 \begin{figure}[H]
  
	

\hspace{+2pt}\includegraphics[width=13.5cm]{Fig_07.jpg}

  
    \caption{Reclassified %MDPI: We revised the hyphen in the figure into minus sign. Please confirm. Please confirm whether the overlapping/incomplete content in this figure affects scientific understanding and if it does, please revise it. % Author's reply: yes, I confirm the revised hyphen in the figure into minus sign. I confirm that the content in this figure is correct and does not affects scientific understanding.
 satellite images after reclassification used by the ‘r.reclass’ technique. Images are from 2019, 2021 and 2023 covering Djibouti. Background: GEBCO. Software: GRASS GIS. Mapping source: Author.}
    \label{fig07}
\end{figure}

Landscape change detection{ was performed }through a comparison among the land cover types detected on the ML-based classified satellite images obtained on different dates. {The} classification of the images is based on an automated evaluation of the spectral reflectance values of the individual pixels on the images, which are assigned to different land cover classes. {The} comparative analysis of {the} classified images {enabled us} to detect landscape changes. {The {observed} transitions {of categories} were the loss of intact vegetation areas; fluctuations in the saline lake, Lake Assal{;} gains in vegetation {and} coastal areas; {and} secondary vegetation growth.} %EE: Please check intended meaning is retained.­ % Author's reply: yes, intended meaning is retained.
These transitions correspond to the predominant patterns of vegetation response to climate warming and growth annual temperatures, which are visible when comparing central and coastal regions of Djibouti, {where the latter have been influenced by the marine climate,} %EE: Please check intended meaning is retained.­ % Author's reply: yes, intended meaning is retained.
especially in the coastal areas of the Gulf of Tadjoura.

{The} traditional (maximal likelihood) and ML{ }classification maps {were} compared {for} 2015, 2019, 2021 and 2023. Comparing the changes on classified images from different dates shows the dynamics of the landscapes over Djibouti from 2015 until 2023. The{ }differences between the{ }land cover or land use types are presented in Figures \ref{fig05}, \ref{fig07} and \ref{fig08}. Figure \ref{fig05} shows the results of the classification made using the maximal likelihood approach; Figure \ref{fig06} shows the accuracy assessment; Figure \ref{fig07} shows the reclassified maps; and Figure \ref{fig08} shows the results of the classification performed using the ML approach. 



The comparison of the classified maps revealed that{ }natural vegetation, such as herbaceous{ }and mosaic shrubland, cropland and grassland, has decreased while the sparse vegetation and bare soil areas typical of desert and semi-desert areas in Djibouti have increased. {Changes} in land cover types in the coastal region of Bab-el-Mandeb and central Djibouti were detected and visualised in the target landscapes over the period from 2015 to 2023. {The} images were compared based on different dates to assess the landscape dynamics, which is defined by the evaluation of land cover types through the identification of various land cover patches, in different regions of Djibouti.

 \begin{figure}[H]
   
\hspace{+2pt}	\includegraphics[width=13.5cm]{Fig_08.jpg}
  
    \caption{Image %MDPI: We revised the hyphen in the figure into minus sign. Please confirm. Please confirm whether the overlapping/incomplete content in this figure affects scientific understanding and if it does, please revise it. % Author's reply: yes, I confirm the revised hyphen in the figure into minus sign. I confirm that the content in this figure is correct and does not affects scientific understanding.
 classification based on the ML gradient boosting classifier algorithm applied for Landsat scenes covering Djibouti in 2015, 2019, 2021 and 2023. Background relief: GEBCO. Software: GRASS GIS. Mapping source: Author.}
    \label{fig08}
\end{figure}

%----------- section ----------->
\section{Conclusions}\label{sec5}

This paper {demonstrated the efficiency of RS data processing using ML methods for visualising} land cover changes. Using this approach, the dynamics of land cover types {has been} detected by{ identifying differences in the} spectral reflectance of pixels {on} images assigned to different categories by ML algorithms of gradient boosting. This method was tested, and has been explained and presented in the form of a series of new maps covering Djibouti. Interdisciplinary problems in geographical sciences{ }often require decisions to be made by diversified approaches{.} ML algorithms, such as the gradient boosting classifier, {improved} mapping workflow {through} programming{ to optimise RS}  data {classification.} When supplementary modules of GIS are used, such tools {point out} the{ }shortcomings of traditional methods in geoinformatics{.} The development of cartographic tools {for RS data processing }using ML {is} a challenging task{.} As technologies mature along with the development of programming algorithms, cartographic modelling and the analysis of landscape changes{ }improve{ by }integrating {applications of scripting methods}. 
Here, we have demonstrated {the use of such} tools using Python's Scikit-Learn library, {adapted} to{ }GRASS GIS. 

Scripting {cartographic} tools{ }enable us to highlight salient aspects of environmental dynamics through the automation of mapping and data visualisation. {They support} geospatial data processing, modelling, visualisation and interpretation. For such geologically complex areas, integrated programming methods assist in handling the issue of data processing in a spatially expressive manner. The{ }effectiveness {of ML }is explained {through} the advanced algorithms of{ }image classification {and automated workflow. }Hence, the use of ML tools is essential for implementing the workflow of processing geographic data to reveal and visualise environmental problems{. }Their applications {demonstrate a} prominent role of{ }scripting for geospatial data handling{ }and{ }environmental analysis. As ML algorithms are developed and implemented in RS and mapping, their performance supports the computational complexity of the tasks involved in cartographic data processing. Many case studies on geographical analysis and mapping, which use conventional GISs, for example, {require highly computational and costly mapping workload} due to diversified cartographic tasks. In such cases, the use of{ }ML, most notably gradient boosting, is an essential solution to optimising cartographic workflow. 

{Future research can use the }thematic maps of land cover types and {continue with the analysis of neighbouring} regions {of the} Bab-el-Mandeb Straight{ }for environmental monitoring. {Moreover, it can} enhance the thematic and topical {research of East Africa, particularly in relation to the geological exploration of Djibouti, environmental analysis, and geophysical monitoring of tectonically active regions within the} %EE: Please check intended meaning is retained.­ % Author's reply: yes, intended meaning is retained.
Afar Triple Junction. {Furthermore, to continue this study, future works can adapt }ML techniques of GRASS GIS{ }for land cover monitoring of {other regions over an extended period }to analyse landscape dynamics {using time series analysis} in regions around the Red Sea, East Africa. %MDPI: The Author Contributions section was removed, as it is unnecessary when there is only one author. Please confirm. % Author's reply: yes, I confirm.
 
\vspace{6pt}
%----------- section ----------->
%\authorcontributions{Supervision, conceptualization, methodology, software, resources, funding acquisition, and project administration, writing---original draft preparation, methodology, software, data curation, visualization, formal analysis, validation, writing---review and editing, and investigation, P.L.}

\funding{The APC was funded by Multidisciplinary Digital Publishing Institute (MDPI).}%MDPI: Please add: ``This research received no external funding'' or ``This research was funded by NAME OF FUNDER grant number XXX.'' and  ``The APC was funded by XXX''. Check carefully that the details given are accurate and use the standard spelling of funding agency names at \url{https://search.crossref.org/funding}, any errors may affect your future funding. % Author's reply: corrected, funding information is added.

\institutionalreview{Not applicable.}

\informedconsent{Not applicable.}

\dataavailability{The data are available in the GitHub repository: %MDPI: Please add the access date (format: Date Month Year), e.g., accessed on 1 January 2020.
\url{https://github.com/paulinelemenkova/GRASS\_ML\_Djibouti}, accessed on 12 June 2025.} 

\acknowledgments{The author thanks the reviewers for reading and reviewing of this manuscript.}

\conflictsofinterest{The author declares no conflicts of interest.} 

\abbreviations{Abbreviations}{
The following abbreviations are used in this manuscript:\\

\noindent 
\begin{tabular}{@{}ll}
{API} & {Application Programming Interface} \\
{CNN} & {Convolutional Neural Network} \\
EO & Earth Observation \\
{FAO} & {Food and Agriculture Organisation} \\
GEBCO & General Bathymetric Chart of the Oceans  \\
GIS & Geographic Information System \\
GMT & Generic Mapping Tool \\
GRASS GIS & Geographic Resources Analysis Support System GIS \\
{KC} & {Kappa Coefficient} \\
{LCCS} & {Land Cover Classification System} \\
{OLI} & {Operational Land Imager} \\
ML & Machine Learning \\
{OA} & {Overall Accuracy} \\
{QGIS} & {Quantum Geographic Information System} \\
{RF} & {Random Forest} \\
RS & Remote Sensing \\
{SVM} & {Support Vector Machine} \\
{TIRS} & {Thermal Infrared Sensor} \\
TIFF & Tag Image File Format \\
USGS & United States Geological Survey \\
UTM & Universal Transverse Mercator \\
WGS84 & World Geodetic System 1984 \\
WRS & Worldwide Reference System \\
\end{tabular}
}

%%%%%%%%%%%%%%%%%%%%%%%%%%%%%%%%%%%%%%%%%%
%% Optional
\appendixtitles{yes} % Leave argument "no" if all appendix headings stay EMPTY (then no dot is printed after "Appendix A"). If the appendix sections contain a heading then change the argument to "yes".
\appendixstart
\appendix

%------------------------ App 1 -------------->
\section[\appendixname~\thesection]{Initial Means for Each Multispectral Band Before Data Processing}\label{AppA}

\vspace{-6pt}
\begin{table}[H]\scriptsize
    \caption{Initial mean values for pixels in spectral bands assigned to target land cover types in Djibouti for the image obtained on 02 July 2015, before iterative classification process.}
  \begin{tabularx}{\textwidth}{p{0.7cm}p{2.5cm}p{1.0cm}p{1.3cm}p{1.0cm}p{1.0cm}p{1.5cm}p{1.5cm}}
  \toprule
  \multirow{2}{*}{\textbf{Class}} & \multicolumn{7}{c}{\textbf{Initial Mean Values for Multispectral Bands of Landsat:  2015}} \\
     & \textbf{1-Coastal Aerosol} & \textbf{2-Blue} & \textbf{3-Green} & \textbf{4-Red} & \textbf{5-NIR} & \textbf{6-SWIR-1} & \textbf{7-SWIR-2} \\
    \midrule
    1 & 8103.97 & 8671.63 & 9562.73 & 9665.43 & 10,%MDPI: We added commas to separate thousands for numbers with five or more digits. Please confirm. The same for the following highlights. % Author's reply: yes, I confirm.
102 & 10,216.9 & 9598.71 \\
    2 & 8421.69 & 9028.36 & 10,061.5 & 10,389.8 & 10,972.8 & 11,250.5 & 10,522.6 \\
    3 & 8739.4 & 9385.08 &10,560.2 & 11,114.2 & 11,843.6 & 12,284.2 & 11,446.6 \\
    4 & 9057.11 & 9741.81 & 11,059 & 11,838.6 & 12,714.5 & 13,317.8 & 12,370.5 \\
    5 & 9374.82 & 10,098.5 & 11,557.7 & 12,563 & 13,585.3 & 14,351.5 & 13,294.4 \\
    6 & 9692.54 & 10,455.3 & 12,056.4 & 13,287.4 & 14,456.2 & 15,385.1 & 14,218.3 \\
    7 & 10,010.2 & 10,812 & 12,555.2 & 14,011.8 & 15,327 & 16,418.8 & 15,142.3 \\
    8 & 10,328 & 11,168.7 & 13,053.9 & 14,736.2 & 16,197.9 & 17,452.4 & 16,066.2 \\
    9 & 10,645.7 & 11,525.4 & 13,552.7 & 15,460.6 & 17,068.7 & 18,486.1 & 16,990.1 \\
    10 & 10,963.4 & 11,882.1 & 14,051.4 & 16,185 & 17,939.6 & 19,519.7 & 17,914 \\
\bottomrule
\end{tabularx}

  \label{tabA1} 
\end{table} 

\vspace{-6pt}
\begin{table}[H]\footnotesize
  \caption{Initial mean values for pixels in spectral bands assigned to target land cover types in Djibouti for the image obtained on 23 July 2017, before iterative classification process.}
  \centering
  \begin{tabularx}{\textwidth}{p{0.7cm}p{2.5cm}p{1.0cm}p{1.3cm}p{1.0cm}p{1.0cm}p{1.5cm}p{1.5cm}}
  \toprule
  \multirow{2}{*}{\textbf{Class}} & \multicolumn{7}{c}{\textbf{Initial Mean Values for Multispectral Bands of Landsat:  2017}} \\
     & \textbf{1-Coastal Aerosol} & \textbf{2-Blue} & \textbf{3-Green} & \textbf{4-Red} & \textbf{5-NIR} & \textbf{6-SWIR-1} & \textbf{7-SWIR-2} \\
    \midrule
    1 & 8074.96 & 8660.92 & 9599.33 & 9725.17 & 10,161.2 & 10,291.9 & 9670.42 \\
    2 & 8390.09 & 9009.56 & 10,079.2 & 10,419.5 & 11,011.6 & 11,300.7 & 10,555.8 \\
    3 & 8705.22 & 9358.2 & 10559.1 & 11,113.9 & 11,861.9 & 12,309.4 & 11441.2 \\
    4 & 9020.35 & 9706.85 & 11,039 & 11,808.3 & 12,712.3 & 13,318.1 & 12,326.7 \\
    5 & 9335.49 & 10,055.5 & 11,518.9 & 12,502.6 & 13,562.7 & 14,326.9 & 13,212.1 \\
    6 & 9650.62 & 10,404.1 & 11,998.8 & 13,197 & 14,413.1 & 15,335.6 & 14,097.5 \\
    7 & 9965.75 &10,752.8 & 12,478.7 &13,891.4 & 15,263.5 & 16,344.3 & 14,982.9 \\
    8 & 10,280.9 & 11,101.4 & 12,958.6 & 14,585.7 & 16,113.8 & 17,353.1 & 15,868.3 \\
    9 & 10,596 &11,450.1 & 13,438.5 & 15,280.1 & 16,964.2 & 18,361.8 & 16,753.7 \\
    10 & 10,911.1 & 11,798.7 & 13,918.4 & 15,974.4 & 17,814.6 & 19,370.5 & 17,639.1 \\
\bottomrule
\end{tabularx}

  \label{tabA2} 
\end{table} 

\vspace{-6pt}
\begin{table}[H]\footnotesize
 \caption{Initial mean values for pixels in spectral bands assigned to target land cover types in Djibouti for the image obtained on 19 August 2021, before iterative classification process.}
  \centering
  \begin{tabularx}{\textwidth}{p{0.7cm}p{2.5cm}p{1.0cm}p{1.3cm}p{1.0cm}p{1.0cm}p{1.5cm}p{1.5cm}}
  \toprule
  \multirow{2}{*}{\textbf{Class}} & \multicolumn{7}{c}{\textbf{Initial Mean Values for Multispectral Bands of Landsat: 2021}} \\
     & \textbf{1-Coastal Aerosol} & \textbf{2-Blue} & \textbf{3-Green} & \textbf{4-Red} &\textbf{ 5-NIR} & \textbf{6-SWIR-1} & \textbf{7-SWIR-2} \\
    \midrule
    1 & 8317.6 & 8806.66 & 9565.01 & 9718.97 & 10,254 & 10,232.6 & 9622.06 \\
    2 & 8578.74 & 9106.8 & 10,005.6 & 10,359.3 & 11,061 & 11,207.5 & 10,480.5 \\
    3 & 8839.87 & 9406.95 & 10,446.1 & 10,999.7 & 11,867.9 & 12,182.4 & 11,338.9 \\
    4 & 9101 & 9707.09 & 10,886.7 & 11,640.1 & 12,674.8 & 13,157.4 & 12,197.4 \\
    5 & 9362.13 &10,007.2 & 11,327.3 & 12,280.4 & 13,481.8 & 14,132.3 & 13,055.8 \\
    6 & 9623.26 & 10,307.4 & 11,767.8 & 12,920.8 & 14,288.7 & 15,107.2 & 13,914.3 \\
    7 & 9884.39 & 10,607.5 & 12,208.4 & 13,561.1 & 15,095.7 & 16,082.1 & 14,772.7 \\
    8 & 10,145.5 & 10,907.7 & 12,649 & 14,201.5 & 15,902.6 & 17,057.1 & 15,631.2 \\
    9 & 10,406.7 & 11,207.8 & 13,089.5 & 14,841.9 & 16,709.6 & 18,032 & 16,489.6 \\
    10 & 10,667.8 & 11,507.9 & 13,530.1 & 15,482.2 & 17,516.5 & 19,006.9 & 17,348 \\
\bottomrule
\end{tabularx}
 
  \label{tabA3} 
\end{table}

\vspace{-6pt}
\begin{table}[H]\footnotesize
  \caption{Initial mean values for pixels in spectral bands assigned to target land cover types in Djibouti for the image obtained on 16 July 2023, before iterative classification process.}
  \centering
  \begin{tabularx}{\textwidth}{p{0.7cm}p{2.5cm}p{1.0cm}p{1.3cm}p{1.0cm}p{1.0cm}p{1.5cm}p{1.5cm}}
  \toprule
  \multirow{2}{*}{\textbf{Class}} & \multicolumn{7}{c}{\textbf{Initial Mean Values for Multispectral Bands of Landsat: 2023}} \\
     & \textbf{1-Coastal Aerosol} & \textbf{2-Blue} & \textbf{3-Green} & \textbf{4-Red} & \textbf{5-NIR} & \textbf{6-SWIR-1} & \textbf{7-SWIR-2} \\
    \midrule
    1 & 8504.75 & 9131.19 & 9993.94 & 10,161.4 & 10,600.6 & 10,554.5 & 9880.58 \\
    2 & 8747.11 & 9404.01 & 10,394.1 & 10,760.9 & 11,378.1 & 11,500.1 & 10,685.6 \\
    3 & 8989.47 & 9676.83 & 10,794.3 & 11,360.4 & 12,155.5 & 12,445.8 & 11,490.6 \\
    4 & 9231.83 & 9949.65 & 11,194.5 & 11,960 & 12,933 & 13,391.4 & 12,295.6 \\
    5 & 9474.19 & 10,222.5 & 11,594.7 & 12,559.5 & 13,710.4 & 14,337 & 13,100.6 \\
    6 & 9716.55 & 10,495.3 & 11,994.9 & 13,159.1 & 14,487.8 & 15,282.7 & 13,905.5 \\
    7 & 9958.91 & 10,768.1 & 12,395.1 & 13,758.6 & 15,265.3 & 16,228.3 & 14,710.5 \\
    8 & 10,201.3 & 11,040.9 & 12,795.3 & 14,358.1 & 16,042.7 &17,173.9 & 15,515.5 \\
    9 & 10,443.6 & 11,313.7 & 13,195.5 & 14,957.7 & 16,820.1 & 18,119.6 & 16,320.5 \\
    10 & 10,686 & 11,586.6 & 13,595.7 & 15,557.2 & 17,597.6 & 19,065.2 & 17,125.5 \\
\bottomrule
\end{tabularx}

  \label{tabA4} 
\end{table}

%------------------------ App 2 -------------->
\section[\appendixname~\thesection]{Iterative Cycles of Image Classification}\label{AppB}

\vspace{-6pt}
\begin{table}[H]\scriptsize  
 \caption{Iterative process of pixel assignment to target classes during image classification: 2015%MDPI: We changed 2021 to 2015, please confirm % Author's reply: yes, I confirm.
.}
  \centering
  \begin{tabularx}{\textwidth}{p{1.6cm}p{1.8cm}p{0.55cm}p{0.55cm}p{0.55cm}p{0.55cm}p{0.55cm}p{0.55cm}p{0.55cm}p{0.55cm}p{0.55cm}p{0.9cm}}
  \toprule
  \multirow{2}{*}{\textbf{2015}} & \multicolumn{10}{c}{\textbf{Pixel Distribution by Categories of Land Cover Classes: 2015}} \\
     & \textbf{Stable Points} %EE: Please check intended meaning is retained.­ % Author's reply: yes, retained.
     & \textbf{1} & \textbf{2} &  \textbf{3} & \textbf{4} & \textbf{5} & \textbf{6} & \textbf{7} & \textbf{8} & \textbf{9} & \textbf{10} \\
    \midrule
    Iteration 1 & 92.24\% & 1672 & 112 &  50 &559       &1110
        &860     &   613   &     390        &635     &   929 \\
    Iteration 2 & 93.78\% & 1568    &    213   &      75     &   584    &   1057&
        850   &     622       & 494   &     647     &   820 \\
    Iteration 3 & 96.28\% & 1538   &     242      &   97        &596     &  1026&
        843    &    658      &  541       & 626     &   763 \\
    Iteration 4 & 96.78\% &1531   &     249     &   117  &      614   &     994&
        840   &     691  &      575    &    599       & 720 \\
    Iteration 5 & 96.97\% & 1531   &     249    &    137    &    622      &  980&
        836    &    714 &       610   &     568      &  683 \\
    Iteration 6 & 97.10\% & 1531  &      249    &    166     &  626   &     963&
        833     &   726  &      643    &    537      &  656 \\
    Iteration 7 & 97.13\% & 1531    &    250    &    194   &     632    &    945&
        834    &    732    &    669     &   512    &    631 \\
    Iteration 8 & 97.11\% & 1531    &    250   &     219     &   650    &   930 &
        827  &      744   &     668    &    498     &   613 \\
    Iteration 9 & 97.52\% & 1531  &      250    &    243    &    658 &      936&
        819     &   732  &      674     &   480     &   607 \\
    Iteration 10 & 97.59\% & 1531   &     250    &    261    &    672   &     941&
        810     &   734    &    659   &     472  &      600 \\
    Iteration 11 & 97.75\% & 1531    &    251  &      279    &    681   &     954&
        798    &    727   &     650  &      462    &    597 \\
    Iteration 12 & 97.79\% & 1531   &     251   &     291   &     704    &    968&
        775   &     720     &   640   &     456     &   594\\
    Iteration 13 & 97.95\% & 1531 &       251   &     305    &    729   &     964 &
        768   &     710    &    629    &    450  &     593 \\
    Iteration 14 & 98.12\% & 1531   &     251    &    314     &   755 &     964&
        761   &     699   &     616    &    448   &     591 \\
\bottomrule
\end{tabularx}
 
  \label{tabB1}
\end{table}

\vspace{-6pt}
\begin{table}[H]\scriptsize  
 \caption{Iterative process of pixel assignment to target classes during image classification: 2017.}
  \centering
\begin{tabularx}{\textwidth}{p{1.6cm}p{1.8cm}p{0.55cm}p{0.55cm}p{0.55cm}p{0.55cm}p{0.55cm}p{0.55cm}p{0.55cm}p{0.55cm}p{0.55cm}p{0.9cm}}
  \toprule
  \multirow{2}{*}{\textbf{2017}} & \multicolumn{10}{c}{\textbf{Pixel Distribution by Categories of Land Cover Classes: 2017}} \\
     & \textbf{Stable Points} %EE: Please check intended meaning is retained.­
     & \textbf{1} & \textbf{2} &  \textbf{3} & \textbf{4} & \textbf{5} & \textbf{6} & \textbf{7} & \textbf{8} &\textbf{ 9} & \textbf{10} \\
    \midrule
    Iteration 1 & 91.34\% & 1698    &    165   &      76    &    574    &   1121 &
        829    &    516    &    427   &     655     &   959 \\
    Iteration 2 & 93.52\% &  1584   &     263    &    117   &     620   &    1047&
        811   &     522   &     531    &    667     &   858 \\
    Iteration 3 & 95.37\% &  1518     &   316   &     137    &    669    &    996&
        788   &     553    &    605   &     640    &    798 \\
    Iteration 4 & 96.70\% & 1486    &    341   &     140     &   699      &  977&
        776  &      578    &    668   &     605     &   750 \\
    Iteration 5 & 97.15\% & 1468    &    357    &   134    &    718     &   972&
        771    &    611    &    707   &     575    &    707 \\
    Iteration 6 & 97.28\% & 1461   &     360  &      131  &      729    &    975&
        783   &     641     &   712    &    561     &   667 \\
    Iteration 7 & 97.51\% & 1456   &     362  &      123    &    741     &  986&
        798     &   651    &    730   &     528    &    645\\
    Iteration 8 & 97.62\% & 1453   &     363  &      118    &    745    &   1016&
        795   &     670    &    725   &     510      &  625 \\
    Iteration 9 & 98.15\% &  1452   &     361 &       120       & 746     &  1031&
        813   &     672 &     714   &     502     &   609\\
\bottomrule
\end{tabularx}
 
  \label{tabB2}
\end{table}

\vspace{-6pt}
\begin{table}[H]\scriptsize  
\caption{Iterative process of pixels' assignment to target classes during image classification: 2021.}
  \centering
\begin{tabularx}{\textwidth}{p{1.6cm}p{1.8cm}p{0.55cm}p{0.55cm}p{0.55cm}p{0.55cm}p{0.55cm}p{0.55cm}p{0.55cm}p{0.55cm}p{0.55cm}p{0.9cm}}
  \toprule
  \multirow{2}{*}{2021} & \multicolumn{10}{c}{\textbf{Pixels' Distribution by Categories of Land Cover Classes: 2021}} \\
     & \textbf{Stable Points} %EE: Please check intended meaning is retained.­
     & \textbf{1} & \textbf{2} &  \textbf{3} & \textbf{4} & \textbf{5} & \textbf{6} & \textbf{7} &\textbf{ 8} & \textbf{9} & \textbf{10} \\
    \midrule
    Iteration 1 & 91.92\% & 1682    &    142    &    128    &    663    &   1062&
        784    &    485     &   371    &    626   &     988 \\
    Iteration 2 & 90.90\% & 1385    &    431    &    174    &    664    &   1017&
        766    &    497    &    467    &    659   &     871 \\
    Iteration 3 & 91.33\% & 1023   &     791     &   193   &     673    &    989&
        758    &    516    &    549      &  634    &    805 \\
    Iteration 4 & 93.64\% & 785  &     1029    &    198   &     684   &     985&
        752   &     540    &    587   &     621     &   750 \\
    Iteration 5 & 95.74\% & 666    &   1147   &     203    &    690   &     984 &
        761  &      557   &     614    &    596    &    713 \\
    Iteration 6 & 96.90\% & 612    &   1201   &     208    &    694   &     984 &
        771      &  573 &       641     &   556     &   691 \\
    Iteration 7 & 97.76\% & 596     &  1217   &     211    &    698     &   985 &
        780    &    593    &    647   &     537    &    667\\
    Iteration 8 & 97.88\% & 589   &    1224   &     216     &   701    &    990 &
        795    &    603   &     645     &   513     &   655 \\
    Iteration 9 & 98.12\% & 588   &    1225   &     223    &    699    &   1007 &
        798  &      608    &    644  &      495   &     644 \\
\bottomrule
\end{tabularx}
  
  \label{tabB3}
\end{table}

\vspace{-6pt}
\begin{table}[H]\footnotesize
 \caption{Iterative process of pixel assignment to target classes during image classification: 2023.}
  \centering
\begin{tabularx}{\textwidth}{p{1.6cm}p{1.8cm}p{0.55cm}p{0.55cm}p{0.55cm}p{0.55cm}p{0.55cm}p{0.55cm}p{0.55cm}p{0.55cm}p{0.55cm}p{0.9cm}}
  \toprule
  \multirow{2}{*}{\textbf{2023}} & \multicolumn{10}{c}{\textbf{Pixel Distribution by Categories of Land Cover classes: 2023}} \\
     & \textbf{Stable Points} %EE: Please check intended meaning is retained.­
     & \textbf{1} & \textbf{2} &  \textbf{3} & \textbf{4} & \textbf{5} & \textbf{6} & \textbf{7} & \textbf{8} &\textbf{ 9} & \textbf{10} \\
    \midrule
    Iteration 1 & 91.95\% & 1665    &     98    &     82    &    574   &   1139 &
        882   &     567    &    370    &    632   &    1032 \\
    Iteration 2 & 91.27\% & 1393    &    367    &    114    &    611   &    1068 &
        871    &    570   &     458    &    689  &      900 \\
    Iteration 3 & 92.44\% & 1134    &    623      &  138     &   641  &     1027 &
        851     &   595    &    538      &  662       &  832 \\
    Iteration 4 & 94.08\% & 939     &   818     &   156   &     653   &    1012 &
        837 &       618      &  588 &        644     &   776 \\
    Iteration 5 & 96.51\% & 872     &   885      &  173      &  661   &    1004 &
        828  &      638    &    609    &    630  &      741 \\
    Iteration 6 & 97.26\% & 847    &    910  &      187   &     668     &  1003 &
        820   &     662 &       606    &    625     &   713 \\
    Iteration 7 & 97.74\% & 842    &    915 &       200    &    676   &     996 &
        829      &  663  &      621   &     606  &      693 \\
    Iteration 8 & 97.77\% & 842   &     915   &     216    &    681&        996 &
        825  &      671      &  635   &     583   &     677 \\
    Iteration 9 & 97.64\% & 842    &    915    &    237     &   683 &       994 &
        829     &   666 &       648      &  569  &      658 \\
    Iteration 10 & 97.88\% & 842    &    915    &    250     &   700     &   983 &
        830    &    666   &     656     &   558    &    641 \\
    Iteration 11 & 97.95\% & 842      &  915  &      264    &    711    &    982 &
        826  &      667  &      655     &   558   &     621 \\
    Iteration 12 & 97.73\% & 842    &    915    &    282   &     719 &        982 &
        822  &      668     &   655    &    557   &     599 \\
    Iteration 13 & 97.74\% & 842   &     915  &      303  &      721    &    980 &
        832     &   664  &      650     &   549    &    585 \\
    Iteration 14 & 97.80\% &  842   &     915   &     332    &    710 &       980 &
        837   &     672     &    639    &    545   &     569 \\
    Iteration 15 & 97.83\% & 842     &   915   &     352    &    723   &     966 &
        844    &    672  &      634     &   533    &    560 \\
    Iteration 16 & 98.13\% & 842   &     915     &   367     &   733 &       965 &
        840    &    672  &       627   &     533    &    547 \\
\bottomrule
\end{tabularx}
 
  \label{tabB4}
\end{table}

%------------------------ App 3 -------------->
\section[\appendixname~\thesection]{Computed Class Means for Land Cover Classes in the Region of Lake Assal, Djibouti, East Africa}\label{AppC}


\centering 
\begin{equation}\label{figSep1}
\scriptsize
\begin{vmatrix}
    & 1 & 2 & 3 & 4 & 5 & 6 & 7 & 8 & 9 & 10 \\
1 & 0 &  &  &  &  &  &  &  &  &  \\
2 & 1.5 & 0 &  &  &  &  &  &  &  &  \\
3 & 5.2 & 3.0 & 0 &  &  &  &  &  &  &  \\
4 & 6.9 & 4.3 & 0.8 & 0 &  &  &  &  &  &  \\
5 & 7.8 & 5.0 & 1.5 & 0.8 & 0 &  &  &  &  &  \\                 
6 & 7.9 & 5.3 & 2.2 & 1.7 & 0.9 & 0 &  &  &  &  \\             
7 & 6.2 & 4.5 & 2.4 & 2.0 & 1.5 & 0.8 & 0 &  &  &  \\                      
8 & 7.4 & 5.5 & 3.4 & 3.1 & 2.6 & 1.8 & 0.8 & 0 &  &  \\
9 & 8.6 & 6.7 & 4.7 & 4.5 & 4.0 & 3.2 & 1.9 & 1.1 & 0 &  \\           
10 & 11.0 & 8.6 & 6.5 & 6.5 & 6.0 & 5.0 & 3.3 & 2.4 & 1.2 & 0 \\                          
\end{vmatrix}
 %
\hspace{3pt}
\begin{vmatrix}
    & 1 & 2 & 3 & 4 & 5 & 6 & 7 & 8 & 9 & 10 \\
1 & 0 &  &  &  &  &  &  &  &  &  \\
2 & 1.3 & 0 &  &  &  &  &  &  &  &  \\
3 & 3.0 & 1.6 & 0 &  &  &  &  &  &  &  \\
4 & 5.8 & 3.5 & 1.2 & 0 &  &  &  &  &  &  \\
5 & 7.1 & 4.4 & 1.8 & 0.8 & 0 &  &  &  &  &  \\                 
6 & 7.4 & 4.9 & 2.3 & 1.6 & 0.9 & 0 &  &  &  &  \\             
7 & 7.5 & 5.2 & 2.9 & 2.4 & 1.8 & 1.0 & 0 &  &  &  \\                      
8 & 7.5 & 5.5 & 3.4 & 3.2 & 2.7 & 2.0 & 1.0 & 0 &  &  \\
9 & 8.0 & 6.1 & 4.2 & 4.2 & 3.8 & 3.1 & 2.2 & 1.2 & 0 &  \\           
10 & 8.8 & 7.1 & 7.1 &  5.2 & 5.3 & 5.0 & 4.3 & 3.4 & 1.1 & 0 \\    
\end{vmatrix} 
\end{equation}
Equation A1: Class separability %MDPI: We have adjusted the format of the figure, please confirm. We recommend changing it to Equation A1 instead of Figure A1 % Author's reply: yes, I confirm the adjusted format. I changed to Equation A1.
 matrices for classified images of Djibouti: 2015 %MDPI: Please confirm if the bold is unnecessary and can be removed. The following highlights are the same. % Author's reply: yes, bold is not necessary and can be removed (I removed it).
 and 2017.

\centering
\begin{equation}\label{figSep2}
\scriptsize
\begin{vmatrix}
    & 1 & 2 & 3 & 4 & 5 & 6 & 7 & 8 & 9 & 10 \\
1 & 0 &  &  &  &  &  &  &  &  &  \\
2 & 1.6 & 0 &  &  &  &  &  &  &  &  \\
3 & 4.1 &   2.5 & 0 &  &  &  &  &  &  &  \\
4 & 5.0 &  3.4 &  0.7 & 0 &  &  &  &  &  &  \\
5 & 6.3 &  4.6  & 1.4 &  0.7 & 0 &  &  &  &  &  \\                 
6 & 6.5 &  4.9 &  2.0  & 1.4 &  0.8 & 0 &  &  &  &  \\             
7 & 5.2  & 4.1 &  2.1 & 1.7 & 1.3 & 0.7 & 0 &  &  &  \\                      
8 & 6.9 & 5.6 & 3.3 & 2.9 & 2.5 & 1.8 & 0.8 & 0 &  &  \\
9 & 5.7 &  4.9 &  3.2  & 2.9 &  2.7 & 2.2  & 1.4 & 0.9 & 0 &  \\           
10 & 8.4 &  7.2 & 5.2 & 4.9 & 4.8 &  4.0  & 2.7 &  2.1 & 0.9 & \\                          
\end{vmatrix}
 %
\hspace{3pt}
\begin{vmatrix}
    & 1 & 2 & 3 & 4 & 5 & 6 & 7 & 8 & 9 & 10 \\
1 & 0 &  &  &  &  &  &  &  &  &  \\
2 & 1.4 & 0 &  &  &  &  &  &  &  &  \\
3 & 4.6  & 3.5 & 0 &  &  &  &  &  &  &  \\
4 & 6.1 &  5.0 &  0.8 & 0 &  &  &  &  &  &  \\
5 & 6.6  & 5.4 & 1.3 & 0.9 & 0 &  &  &  &  &  \\                 
6 & 6.8 &  5.9 & 1.9 & 1.5 & 0.8 & 0 &  &  &  &  \\             
7 & 6.8 & 5.9 & 2.5 & 2.3 & 1.7 & 1.0 & 0 &  &  &  \\                      
8 & 5.5 & 4.8 & 2.6 &  2.5 & 2.0 & 1.5 & 0.8 & 0 &  &  \\
9 & 7.6 & 6.8 & 4.2 & 4.1&  3.6 & 3.0 & 1.9 & 0.8 & 0 &  \\           
10 & 9.0 & 8.2 & 5.4 & 5.5 & 4.9 & 4.2 & 3.1 & 1.7 & 1.0 & 0 \\    
\end{vmatrix} 
\end{equation}
Equation A2: Class %MDPI: We have adjusted the format of the figure, please confirm. We recommend changing it to Equation A2 instead of Figure A2. % Author's reply: yes, I confirm the adjusted format. I changed to Equation A2.
 separability matrices for classified images of Djibouti: 2021 and 2023.

%------------------------ App 4 -------------->
\section[\appendixname~\thesection]{Computed Class Means for Land Cover Classes in the Region of Lake Assal, Djibouti, East Africa}\label{AppD}

\vspace{-6pt}
\begin{table}[H]\small
\caption{Computed %MDPI: Please confirm the alignment change. % Author's reply: yes, I confirm and agree.
 class means of 10 land cover types in Lake Assal, Djibouti. The computations are based on Digital Numbers (DNs) of pixels in the evaluated landscapes. Data: multispectral Landsat 8--9 OLI/TIRS images, band (1 to 7) for years 2015 to 2023.} \label{tabAppD}

\begin{tabularx}{\textwidth}{CCCCCCCC}
\toprule \textbf{Class} & \textbf{Band 1} & \textbf{Band 2} & \textbf{Band 3} & \textbf{Band 4} & \textbf{Band 5} & \textbf{Band 6} & \textbf{Band 7} \\ \midrule 


2015 &  &  &  &  &  &  & \\
\midrule
1 & 7394.23 & 7878.92 & 8615.8 & 8244.47 & 8385.55 & 8786 & 8586.51 \\
2 & 8998.3 & 9445.8 & 9981.23 & 9575.59 & 9422.92 & 9284.88 & 8887.37 \\
3 & 9230.46 & 9913.19 & 11,%MDPI: We added commas to separate thousands for numbers with five or more digits. Please confirm. The same for the following highlights. % Author's reply: yes, I confirm.
138.2 & 11,958.2 & 12,864.5 & 12,862.9 & 11,816.5 \\
4 & 9430.8 & 10,146.6 & 11,465.7 & 12,446 & 13,497.2 & 13,806.2 & 12,619 \\
5 & 9547.38 & 10,303 & 11,752.4 & 12,912.6 & 14,092.7 & 14,641.3 & 13,360.3 \\
6 & 9821.76 & 10,624.5 & 12,235.6 & 13,635.1 & 14,896.9 & 15,478.6 & 14,133.9 \\
7 & 10,246.8 & 11,081.6 & 12,862.4 & 14,552 & 16,012.4 & 16,653.6 & 15,224.3 \\
8 & 10,661.7 & 11,544 & 13,560.5 & 15,543.5 & 17,259.8 & 18,626.3 & 17,001.6 \\
9 & 11,251.3 & 12,246.2 & 14,758.7 & 17,309 & 19,326.1 & 21,064.2 & 19,184.3 \\
10 & 11,881.4 & 12,946.1 & 15,865.1 & 18,729.9 & 20,856 & 24,277.4 & 22,627.1 \\
\midrule
2017 &  &  &  & & & & \\
\midrule
1 & 7371.5 & 7875.98 & 8664.69 & 8356.38 & 8565.28 & 8995.22 & 8778.1 \\
2 & 8858.16 & 9357.22 & 9951.7 & 9564.73 & 9444.54 & 9371.12 & 8992.18 \\
3 & 8945.89 & 9777.33 & 11,054.5 & 11,333.4 & 11,764 & 11,571.8 & 10,743 \\
4 & 9321.25 & 10,018.9 & 11,324 & 12,242.3 & 13,244.4 & 13,468.5 & 12,268.8 \\
5 & 9496.17 & 10,212.7 & 11,611.2 & 12,696.3 & 13,855.2 & 14,375.8 & 13,058.1 \\
6 & 9715.96 & 10,479.6 & 12,049.1 & 13,373.4 & 14,635.2 & 15,185.6 & 13,795 \\
7 & 10,087 & 10,877.3 & 12,605.2 & 14,223.3 & 15,670.3 & 16,380 & 14,930.2 \\
8 & 10,526.8 & 11378.8 & 13341.7 & 15248 & 16971.7 & 18216.1 & 16582.3 \\
9 & 11,127.8 & 12,082.7 & 14,519.5 & 16,951.3 & 19,046.9 & 20,773.9 & 18,771.8 \\
10 & 11,728.2 & 12,813.2 & 15,676 & 18,406.1 & 20,794.3 & 24,015.7 & 22,082.7 \\



\bottomrule
\end{tabularx}
\end{table}

\begin{table}[H]\ContinuedFloat
\small
\caption{{\em Cont.}}

\begin{tabularx}{\textwidth}{CCCCCCCC}
\toprule \textbf{Class} & \textbf{Band 1} & \textbf{Band 2} & \textbf{Band 3} & \textbf{Band 4} & \textbf{Band 5} & \textbf{Band 6} & \textbf{Band 7} \\ \midrule 


2021 &  &  &  & & & & \\
\midrule
1 & 7281.64 & 7688.52 & 8323.58 & 8019.43 & 8124.06 & 8480.67 & 8350.31 \\
2 & 8578.45 & 8990.99 & 9497.29 & 9305.59 & 9360.4 & 9422.46 & 9107.76 \\
3 & 9032.4 & 9610.46 & 10,681.6 & 11,350.9 & 12,512.9 & 12,392.4 & 11,404.3 \\
4 & 9232.84 & 9855.5 & 11063 & 11,913.4 & 13,238.9 & 13,353.2 & 12,168.5 \\
5 & 9378.87 & 10,036.2 & 11,347.7 & 12,369.6 & 13,788.5 & 14,219 & 12,929 \\
6 & 9563.91 & 10,256.7 & 11,715.6 & 12,988.2 & 14,561.6 & 15,110.5 & 13,758.7 \\
7 & 10,037.5 & 10,785.2 &12,429.3 & 13,969.7 & 15,629.4 & 16,072.7 & 14,651.6 \\
8 & 10,332.3 & 11,120.3 & 12,974.7 & 14,809.1 & 16,765.2 & 17,972.9 & 16,439.1 \\
9 & 10,934.6 & 11,846.4 & 14,132.9 & 16,407.4 & 18,587 & 20,167.3 & 18,278.5 \\
10 & 11,319.5 & 12,358.9 & 15,074.5 & 17,683.4 & 20,102 & 23,286.8 & 21,439.7 \\
\midrule
2023 &  &  &  & & & & \\
\midrule
1 & 7473.59 & 8079.42 & 8857.77 & 8532.88 & 8642.03 & 8940.82 & 8737.73 \\
2 & 9179.37 & 9560.93 & 9889.56 & 9619.59 & 9502.14 & 9260.89 & 8893.31 \\
3 & 9285.35 & 10,023 & 11,304.4 & 12,097.2 & 13,066.2 & 13,055 & 11,870.5 \\
4 & 9260.65 & 9993.24 & 11324.5 & 12,294.7 & 13,624 & 14,111.4 & 12,725.8 \\
5 & 9612.99 & 10,407.7 & 11,850 & 12,914.8 & 14,185.6 & 14,445.7 & 13,030.6 \\
6 & 9612.99 & 10,407.7 & 11,850 & 12,914.8 & 14,185.6 & 14,445.7 & 13,030.6 \\
7 & 10,105.5 & 10,948.7 & 12,646.3 & 14,221.3 & 15,852.4 & 16,551.5 & 14,997.2 \\
8 & 10,551 & 11,459.3 & 13,401.9 & 15,287.1 & 17,177.8 & 18,217.3 & 16,436.4 \\
9 & 10,830.2 & 11,812.4 & 14,092.6 & 16,381.5 & 18,713.8 & 20,687.4 & 18,410.8 \\
10 & 11,191.6 & 12,288.7 & 14,890.2 & 17,469.4 & 20,034.6 & 23,154.3 & 20,893.2 \\
\bottomrule
\end{tabularx}
\end{table}

%%%%%%%%%%%%%%%%%%%%%%%%%%%%%%%%%%%%%%%%%%
\begin{adjustwidth}{-\extralength}{0cm}
%\printendnotes[custom] % Un-comment to print a list of endnotes

\reftitle{References}
%\nocite{*}

%=====================================
% References, variant A: external bibliography
%=====================================
\begin{thebibliography}{999}

\bibitem[Ndehedehe(2022)]{Ndehedehe2022}
Ndehedehe, C. Remote Sensing of Surface Vegetation.
\newblock In {\em Satellite Remote Sensing of Terrestrial Hydrology}; Springer
  International Publishing: Cham, Switzerland, %MDPI: Newly added information, please confirm. The same for the following highlights. % Author's reply: yes, I confirm all changes in the References.
 2022; pp. 131--176.
\newblock {\url{https://doi.org/10.1007/978-3-030-99577-5_7}}.

\bibitem[Foli et~al.(2022)Foli, Williams, Boakye, Azumah, Agyekum, and
  Wiafe]{Foli}
Foli, B.A.K.; Williams, I.K.; Boakye, A.A.; Azumah, D.M.Y.; Agyekum, K.A.;
  Wiafe, G.
\newblock Earth Observation Services in Support of West Africa's Blue Economy:
  Coastal Resilience and Climate Impacts.
\newblock {\em Remote Sens. Earth Syst. Sci.} {\bf 2022}, {\em
  5},~59--70.
\newblock {\url{https://doi.org/10.1007/s41976-021-00058-x}}.

\bibitem[Ifejika~Speranza et~al.(2023)Ifejika~Speranza, Akinyemi, Baratoux,
  Benveniste, Ceperley, Driouech, and Helmschrot]{Speranza}
Ifejika~Speranza, C.; Akinyemi, F.O.; Baratoux, D.; Benveniste, J.; Ceperley,
  N.; Driouech, F.; Helmschrot, J.
\newblock Enhancing the Uptake of Earth Observation Products and Services in
  Africa Through a Multi-level Transdisciplinary Approach.
\newblock {\em Surv. Geophys.} {\bf 2023}, {\em 44},~7--41.
\newblock {\url{https://doi.org/10.1007/s10712-022-09724-1}}.

\bibitem[Springer et~al.(2023)Springer, Lopez, Owor, Frappart, and
  Stieglitz]{Springer2023}
Springer, A.; Lopez, T.; Owor, M.; Frappart, F.; Stieglitz, T.
\newblock The Role of Space-Based Observations for Groundwater Resource
  Monitoring over Africa.
\newblock {\em Surv. Geophys.} {\bf 2023}, {\em 44},~123--172.
\newblock {\url{https://doi.org/10.1007/s10712-022-09759-4}}.

\bibitem[Jia et~al.(2023)Jia, Li, Viña, Cheng, Dou, Song, He, and
  Liu]{10281679}
Jia, N.; Li, Y.; Viña, A.; Cheng, J.; Dou, Y.; Song, Q.; He, L.; Liu, J.
\newblock Long Image Time Series for Crop Extraction Based on the Automatically
  Generated Samples Algorithm.
\newblock In Proceedings of the IGARSS 2023-2023 IEEE International
  Geoscience and Remote Sensing Symposium,  Pasadena, CA, USA, 16--21 July 2023;  pp. 3502--3505.
\newblock {\url{https://doi.org/10.1109/IGARSS52108.2023.10281679}}.

\bibitem[Ngondo et~al.(2021)Ngondo, Mango, Liu, Nobert, Dubi, and
  Cheng]{su13084092}
Ngondo, J.; Mango, J.; Liu, R.; Nobert, J.; Dubi, A.; Cheng, H.
\newblock Land-Use and Land-Cover (LULC) Change Detection and the Implications
  for Coastal Water Resource Management in the Wami–Ruvu Basin, Tanzania.
\newblock {\em Sustainability} {\bf 2021}, {\em 13}, 4092.
\newblock {\url{https://doi.org/10.3390/su13084092}}.

\bibitem[Natarajan et~al.(2016)Natarajan, Latva-Käyrä, Zyadin, and
  Pelkonen]{NATARAJAN2016256}
Natarajan, K.; Latva-Käyrä, P.; Zyadin, A.; Pelkonen, P.
\newblock {New methodological approach for biomass resource assessment in India
  using GIS application and land use/land cover (LULC) maps}.
\newblock {\em Renew. Sustain. Energy Rev.} {\bf 2016}, {\em
  63},~256--268.
\newblock {\url{https://doi.org/10.1016/j.rser.2016.05.070}}.

\bibitem[Borana and Yadav(2023)]{BORANA2023261}
Borana, S.; Yadav, S.
\newblock Chapter 10-Urban land-use susceptibility and sustainability---Case
  study. In {\em Water, Land, and Forest Susceptibility and Sustainability};
  Chatterjee, U., Pradhan, B., Kumar, S., Saha, S., Zakwan, M., Fath, B.D.,
  Fiscus, D., Eds.; Academic Press: Cambridge, MA, USA, 
  2023; Volume~2, %MDPI: We removed ``{\em Science of Sustainable   Systems}'', please confirm % Author's reply: yes, I confirm.
 pp. 261--286.
\newblock {\url{https://doi.org/10.1016/B978-0-443-15847-6.00010-0}}.

\bibitem[Zhu et~al.(2021)Zhu, Zhao, Liang, and Qin]{Zhu02012021}
Zhu, A.X.; Zhao, F.H.; Liang, P.; Qin, C.Z.
\newblock Next generation of GIS: must be easy.
\newblock {\em Ann.  GIS} {\bf 2021}, {\em 27},~71--86.
\newblock {\url{https://doi.org/10.1080/19475683.2020.1766563}}.

\bibitem[Dyson et~al.(2024)Dyson, Nicolau, Saah, and Clinton]{Dyson2024}
Dyson, K.; Nicolau, A.P.; Saah, D.; Clinton, N. Image Manipulation: Bands,
  Arithmetic, Thresholds, and Masks.
\newblock In {\em Cloud-Based Remote Sensing with Google Earth Engine:
  Fundamentals and Applications}; Springer International Publishing: Cham, Switzerland,
  2024; pp. 97--114.
\newblock {\url{https://doi.org/10.1007/978-3-031-26588-4_5}}.

\bibitem[Kochher and Sharma(2016)]{7877513}
Kochher, R.; Sharma, A.
\newblock {Improved principle component analysis based gray stretch algorithm
  for landsat image segmentation}.
\newblock In Proceedings of the 2016 2nd International Conference on Next
  Generation Computing Technologies (NGCT), Dehradun, India, 14--16 October 2016; pp. 765--771.
\newblock {\url{https://doi.org/10.1109/NGCT.2016.7877513}}.

\bibitem[Ferchichi et~al.(2024)Ferchichi, Chihaoui, and
  Ferchichi]{FERCHICHI2024122211}
Ferchichi, A.; Chihaoui, M.; Ferchichi, A.
\newblock Spatio-temporal modeling of climate change impacts on drought
  forecast using Generative Adversarial Network: A case study in Africa.
\newblock {\em Expert Syst. Appl.} {\bf 2024}, {\em 238},~122211.
\newblock {\url{https://doi.org/10.1016/j.eswa.2023.122211}}.

\bibitem[Balha and Singh(2022)]{Balha2022}
Balha, A.; Singh, C.K. Comparison of Maximum Likelihood, Neural Networks, and
  Random Forests Algorithms in Classifying Urban Landscape.
\newblock In {\em Application of Remote Sensing and GIS in Natural Resources
  and Built Infrastructure Management}; Springer International Publishing:
  Cham, Switzerland, 2022; pp. 29--38.
\newblock {\url{https://doi.org/10.1007/978-3-031-14096-9_2}}.

\bibitem[Lemenkova(2024)]{Lemenkova202401}
Lemenkova, P.
\newblock {Random Forest Classifier Algorithm of Geographic Resources Analysis
  Support System Geographic Information System for Satellite Image Processing:
  Case Study of Bight of Sofala, Mozambique}.
\newblock {\em Coasts} {\bf 2024}, {\em 4},~127--149.
\newblock {\url{https://doi.org/10.3390/coasts4010008}}.

\bibitem[Zhang et~al.(2023)Zhang, Zhao, Xue, Meng, and Wang]{10450641}
Zhang, D.; Zhao, T.; Xue, S.; Meng, Y.; Wang, W.
\newblock Diagnosis of Dark Field Partial Slight Discharge of GIS Coaxial
  Busbar Using KPCA-SVM-Based Heuristic Sequence Comparison Method.
\newblock In Proceedings of the 2023 China Automation Congress (CAC),  Chongqing, China, 17--19 November 2023;
  pp. 4487--4492.
\newblock {\url{https://doi.org/10.1109/CAC59555.2023.10450641}}.

\bibitem[Suárez-Fernández et~al.(2024)Suárez-Fernández, Martínez-Sánchez,
  and Arias]{SUAREZFERNANDEZ2024}
Suárez-Fernández, G.E.; Martínez-Sánchez, J.; Arias, P.
\newblock Assessment of vegetation indices for mapping burned areas using a
  deep learning method and a comprehensive forest fire dataset from Landsat
  collection.
\newblock {\em Adv. Space Res.} {\bf 2025%MDPI: We changed year, please confirm % Author's reply: yes, I confirm.
}, \emph{75}, 1665--1685.
\newblock {\url{https://doi.org/10.1016/j.asr.2024.12.001}}.

\bibitem[Huang et~al.(2020)Huang, Cao, Guo, Jiang, Li, and
  Guo]{HUANG2020104580}
Huang, F.; Cao, Z.; Guo, J.; Jiang, S.H.; Li, S.; Guo, Z.
\newblock Comparisons of heuristic, general statistical and machine learning
  models for landslide susceptibility prediction and mapping.
\newblock {\em Catena} {\bf 2020}, {\em 191},~104580.
\newblock {\url{https://doi.org/10.1016/j.catena.2020.104580}}.

\bibitem[Arabameri et~al.(2021)Arabameri, Arora, Pal, Mitra, Saha, Nalivan,
  Panahi, and Moayedi]{Arabameri}
Arabameri, A.; Arora, A.; Pal, S.C.; Mitra, S.; Saha, A.; Nalivan, O.A.;
  Panahi, S.; Moayedi, H.
\newblock K-Fold and State-of-the-Art Metaheuristic Machine Learning Approaches
  for Groundwater Potential Modelling.
\newblock {\em Water Resour. Manag.} {\bf 2021}, {\em 35},~1837--1869.
\newblock {\url{https://doi.org/10.1007/s11269-021-02815-5}}.

\bibitem[Hara~Sudhan et~al.(2018)Hara~Sudhan, Aravind, and Thanaraj]{8524429}
Hara~Sudhan, G.H.; Aravind, R.G.; Thanaraj, K.P.
\newblock Multispectral Analysis of Satellite Images Using Heuristic Algorithm.
\newblock In Proceedings of the 2018 International Conference on Communication
  and Signal Processing (ICCSP),  Chennai, India, 3--5 April 2018; pp. 1010--1013.
\newblock {\url{https://doi.org/10.1109/ICCSP.2018.8524429}}.

\bibitem[Vitale et~al.(2024)Vitale, Salvo, and Lamonaca]{10615674}
Vitale, A.; Salvo, C.; Lamonaca, F.
\newblock A Novel Geospatial Methodology for Measuring and Mapping
  Spatiotemporal Built-Up Dynamics Based on Google Earth Engine and
  Unsupervised K-Means Clustering of Multispectral Satellite Imagery.
\newblock In Proceedings of the 2024 IEEE International Workshop on Metrology
  for Living Environment (MetroLivEnv),  Chania, Greece, 12--14 June 2024; pp. 57--62.
\newblock {\url{https://doi.org/10.1109/MetroLivEnv60384.2024.10615674}}.

\bibitem[Kirkby et~al.(1997)Kirkby, Badcock, and Eklund]{669164}
Kirkby, S.; Badcock, J.; Eklund, P.
\newblock Encoding heuristic knowledge for GIS classification.
\newblock In Proceedings of the 1997 IEEE International Conference on
  Intelligent Processing Systems (Cat. No.97TH8335),  Beijing, China,  28--31 October 1997; Volume~2, pp.
  1143--1145.
\newblock {\url{https://doi.org/10.1109/ICIPS.1997.669164}}.

\bibitem[Hong and Bian(2008)]{4666175}
Hong, Z.; Bian, F.
\newblock A Heuristic Approach for Fast Mining Association Rules in
  Transportation System.
\newblock In Proceedings of the 2008 Fifth International Conference on Fuzzy
  Systems and Knowledge Discovery,  Jinan, China, 18--20 October 2008; Volume~2, pp. 541--545.
\newblock {\url{https://doi.org/10.1109/FSKD.2008.332}}.

\bibitem[Pham et~al.(2020)Pham, Nguyen, Nguyen, Pham, Vu, and Bui]{8999627}
Pham, V.D.; Nguyen, Q.H.; Nguyen, H.D.; Pham, V.M.; Vu, V.M.; Bui, Q.T.
\newblock Convolutional Neural Network—Optimized Moth Flame Algorithm for
  Shallow Landslide Susceptible Analysis.
\newblock {\em IEEE Access} {\bf 2020}, {\em 8},~32727--32736.
\newblock {\url{https://doi.org/10.1109/ACCESS.2020.2973415}}.

\bibitem[Han et~al.(2021)Han, Liao, Gao, Wang, and Yang]{9708387}
Han, S.; Liao, S.; Gao, F.; Wang, B.; Yang, N.
\newblock Pattern Recognition of UHF Partial Discharge Signals in GIS Based on
  Deep Learning.
\newblock In Proceedings of the 2021 International Conference on Advanced
  Electrical Equipment and Reliable Operation (AEERO),   Beijing, China, 15--17 October 2021; pp. 1--5.
\newblock {\url{https://doi.org/10.1109/AEERO52475.2021.9708387}}.

\bibitem[Badr and Oommen(2006)]{1634653}
Badr, G.; Oommen, B.
\newblock On optimizing syntactic pattern recognition using tries and AI-based
  heuristic-search strategies.
\newblock {\em IEEE Trans. Syst. Man, Cybern. Part B
  Cybernetics} {\bf 2006}, {\em 36},~611--622.
\newblock {\url{https://doi.org/10.1109/TSMCB.2005.861860}}.

\bibitem[Yang and Zheng(2023)]{10145147}
Yang, J.; Zheng, X.
\newblock Partial Discharge Pattern Recognition in GIS Based on Multiscale
  Dispersion Entropy and Stacking Ensemble Learning.
\newblock In Proceedings of the 2023 International Conference on Pattern
  Recognition, Machine Vision and Intelligent Algorithms (PRMVIA),  Beihai, China, 24--26 March 2023; pp.
  81--85.
\newblock {\url{https://doi.org/10.1109/PRMVIA58252.2023.00020}}.

\bibitem[Sasmi~Hidayatul et~al.(2019)Sasmi~Hidayatul, Djunaidy, and
  Muklason]{8938484}
Sasmi~Hidayatul, Y.T.; Djunaidy, A.; Muklason, A.
\newblock Solving Multi-objective Vehicle Routing Problem Using Hyper-heuristic
  Method By Considering Balance of Route Distances.
\newblock In Proceedings of the 2019 International Conference on Information
  and Communications Technology (ICOIACT),  Yogyakarta, Indonesia,  24--25 July 2019; pp. 937--942.
\newblock {\url{https://doi.org/10.1109/ICOIACT46704.2019.8938484}}.

\bibitem[Tang and Cao(2020)]{9279780}
Tang, Z.; Cao, Z.
\newblock Application of Convolutional Neural Network Transfer Learning in
  Partial Discharge Pattern Recognition.
\newblock In Proceedings of the 2020 IEEE International Conference on High
  Voltage Engineering and Application (ICHVE),  Beijing, China, 6--10 September 2020; pp. 1--4.
\newblock {\url{https://doi.org/10.1109/ICHVE49031.2020.9279780}}.

\bibitem[Neteler et~al.(2012)Neteler, Bowman, Landa, and Metz]{NETELER2012124}
Neteler, M.; Bowman, M.H.; Landa, M.; Metz, M.
\newblock GRASS GIS: A multi-purpose open source GIS.
\newblock {\em Environ. Model. Softw.} {\bf 2012}, {\em
  31},~124--130.
\newblock {\url{https://doi.org/10.1016/j.envsoft.2011.11.014}}.

\bibitem[Team et~al.(2025)Team, Landa, Neteler, Metz, Petr{\'a}{\v s}ov{\'a},
  Petr{\'a}{\v s}, Clements, Zigo, Larsson, Kladivov{\'a}, Haedrich,
  Blumentrath, Andreo, Cho, Gebbert, Narti{\v s}s, Kudrnovsky, Delucchi,
  Zambelli, Lennert, Mit{\'a}{\v s}ov{\'a}, Chemin, Pe{\v s}ek, Barton,
  Tawalika, Ovsienko, and Bowman]{grass_development_team_2025_14918754}
Team, G.D.; Landa, M.; Neteler, M.; Metz, M.; Petr{\'a}{\v s}ov{\'a}, A.;
  Petr{\'a}{\v s}, V.; Clements, G.; Zigo, T.; Larsson, N.; Kladivov{\'a}, L.;
  et~al.
\newblock GRASS GIS, 2025. Online resource. URL of the website: \url{https://doi.org/10.5281/zenodo.14918754}. Accessed 12 June 2025. %MDPI: Please provide more information about the article type, such as book (please provide the name and location of the publisher); online resource (please provide the URL of the website and the date it was accessed (Date Month Year)); or journal article (please provide the name of the journal, the year and volume in which it was published, and the page number). Please refer to https://www.mdpi.com/authors/references for full reference formatting guides. % Author's reply:  corrected, new information is added.

\bibitem[McCabe(2019)]{McCabe02102019}
McCabe, R.
\newblock Policing the Seas: Building Constabulary Maritime Governance in the
  Horn of Africa–The Case of Djibouti and Kenya.
\newblock {\em Afr. Secur.} {\bf 2019}, {\em 12},~330--355.
\newblock {\url{https://doi.org/10.1080/19392206.2019.1667053}}.

\bibitem[Lawale and Ahmad(2023)]{Lawale02102023}
Lawale, S.; Ahmad, T.
\newblock The Role of Small States in Power Contestations in the Horn of Africa
  Case Study of Djibouti.
\newblock {\em Asian J. Middle East. Islam. Stud.} {\bf 2023},
  {\em 17},~325--340.
\newblock {\url{https://doi.org/10.1080/25765949.2023.2300581}}.

\bibitem[Verdugo(2016)]{Verdugo2016}
Verdugo, D. Djibouti.
\newblock In {\em Encyclopedia of Tourism}; Springer International Publishing:
  Cham,  Switzerland, 2016; pp. 266--267.
\newblock {\url{https://doi.org/10.1007/978-3-319-01384-8_692}}.

\bibitem[Haile et~al.(2020)Haile, Tang, Hosseini-Moghari, Liu, Gebremicael,
  Leng, Kebede, Xu, and Yun]{Haile}
Haile, G.G.; Tang, Q.; Hosseini-Moghari, S.M.; Liu, X.; Gebremicael, T.G.;
  Leng, G.; Kebede, A.; Xu, X.; Yun, X.
\newblock Projected Impacts of Climate Change on Drought Patterns Over East
  Africa.
\newblock {\em Earth's Future} {\bf 2020}, {\em 8},~e2020EF001502.
\newblock {\url{https://doi.org/10.1029/2020EF001502}}.

\bibitem[{Gebremeskel Haile} et~al.(2020){Gebremeskel Haile}, Tang, Leng, Jia,
  Wang, Cai, Sun, Baniya, and Zhang]{GEBREMESKELHAILE2020135299}
{Gebremeskel Haile}, G.; Tang, Q.; Leng, G.; Jia, G.; Wang, J.; Cai, D.; Sun,
  S.; Baniya, B.; Zhang, Q.
\newblock Long-term spatiotemporal variation of drought patterns over the
  Greater Horn of Africa.
\newblock {\em Sci. Total. Environ.} {\bf 2020}, {\em 704},~135299.
\newblock {\url{https://doi.org/10.1016/j.scitotenv.2019.135299}}.

\bibitem[Lemenkova(2022)]{Lemenkova202221}
Lemenkova, P.
\newblock {Evapotranspiration, vapour pressure and climatic water deficit in
  Ethiopia mapped using GMT and TerraClimate dataset}.
\newblock {\em J. Water Land Dev.} {\bf 2022}, {\em
  54},~201--209.
\newblock {\url{https://doi.org/10.24425/jwld.2022.141573}}.

\bibitem[AghaKouchak(2015)]{AGHAKOUCHAK2015127}
AghaKouchak, A.
\newblock A multivariate approach for persistence-based drought prediction:
  Application to the 2010–2011 East Africa drought.
\newblock {\em J. Hydrol.} {\bf 2015}, {\em 526},~127--135. %MDPI: We removed ``Drought processes, modeling, and mitigation,'', please confirm % Author's reply: yes, I confirm.
\newblock 
  {\url{https://doi.org/10.1016/j.jhydrol.2014.09.063}}.

\bibitem[Agutu et~al.(2017)Agutu, Awange, Zerihun, Ndehedehe, Kuhn, and
  Fukuda]{AGUTU2017287}
Agutu, N.; Awange, J.; Zerihun, A.; Ndehedehe, C.; Kuhn, M.; Fukuda, Y.
\newblock Assessing multi-satellite remote sensing, reanalysis, and land
  surface models' products in characterizing agricultural drought in East
  Africa.
\newblock {\em Remote Sens. Environ.} {\bf 2017}, {\em 194},~287--302.
\newblock {\url{https://doi.org/10.1016/j.rse.2017.03.041}}.

\bibitem[D’Ercole et~al.(2024)D’Ercole, Casella, Panegrossi, and
  Sanò]{DERCOLE2024104264}
D’Ercole, R.; Casella, D.; Panegrossi, G.; Sanò, P.
\newblock A high temporal resolution NDVI time series to monitor drought events
  in the Horn of Africa.
\newblock {\em Int. J. Appl. Earth Obs. Geoinf.} {\bf 2024}, {\em 135},~104264.
\newblock {\url{https://doi.org/10.1016/j.jag.2024.104264}}.

\bibitem[Tufa~Dinku and Connor(2011)]{Dinku10112011}
Tufa~Dinku, P.C.; Connor, S.J.
\newblock Challenges of satellite rainfall estimation over mountainous and arid
  parts of east Africa.
\newblock {\em Int. J. Remote Sens.} {\bf 2011}, {\em
  32},~5965--5979.
\newblock {\url{https://doi.org/10.1080/01431161.2010.499381}}.

\bibitem[Mishra and Singh(2010)]{MISHRA2010202}
Mishra, A.K.; Singh, V.P.
\newblock {A Review of Drought Concepts}.
\newblock {\em J. Hydrol.} {\bf 2010}, {\em 391},~202--216.
\newblock {\url{https://doi.org/10.1016/j.jhydrol.2010.07.012}}.

\bibitem[Fava and Vrieling(2021)]{FAVA202144}
Fava, F.; Vrieling, A.
\newblock Earth observation for drought risk financing in pastoral systems of
  sub-Saharan Africa.
\newblock {\em Curr. Opin. Environ. Sustain.} {\bf 2021},
  {\em 48},~44--52. %MDPI: We removed ``The dryland social-ecological systems in changing environments,'', please confirm % Author's reply: yes, I confirm.
\newblock 
  {\url{https://doi.org/10.1016/j.cosust.2020.09.006}}.

\bibitem[Bravar and Kavvas(1991)]{BRAVAR1991281}
Bravar, L.; Kavvas, M.
\newblock {On the Physics of Droughts. I. A Conceptual Framework}.
\newblock {\em J. Hydrol.} {\bf 1991}, {\em 129},~281--297.
\newblock {\url{https://doi.org/10.1016/0022-1694(91)90055-M}}.

\bibitem[Schulman(2019)]{Schulman02012019}
Schulman, S.
\newblock Climate Change Challenges and Djibouti.
\newblock {\em  RUSI J.} {\bf 2019}, {\em 164},~62--75.
\newblock {\url{https://doi.org/10.1080/03071847.2019.1605020}}.

\bibitem[Carrão et~al.(2016)Carrão, Naumann, and Barbosa]{CARRAO2016108}
Carrão, H.; Naumann, G.; Barbosa, P.
\newblock {Mapping Global Patterns of Drought Risk: An Empirical Framework
  Based on Sub-National Estimates of Hazard, Exposure and Vulnerability}.
\newblock {\em Glob. Environ. Change} {\bf 2016}, {\em 39},~108--124.
\newblock {\url{https://doi.org/10.1016/j.gloenvcha.2016.04.012}}.

\bibitem[Fenta et~al.(2020)Fenta, Tsunekawa, Haregeweyn, Poesen, Tsubo,
  Borrelli, Panagos, Vanmaercke, Broeckx, Yasuda, Kawai, and
  Kurosaki]{FENTA2020135016}
Fenta, A.A.; Tsunekawa, A.; Haregeweyn, N.; Poesen, J.; Tsubo, M.; Borrelli,
  P.; Panagos, P.; Vanmaercke, M.; Broeckx, J.; Yasuda, H.;  et~al.
\newblock Land susceptibility to water and wind erosion risks in the East
  Africa region.
\newblock {\em Sci. Total. Environ.} {\bf 2020}, {\em 703},~135016.
\newblock {\url{https://doi.org/10.1016/j.scitotenv.2019.135016}}.

\bibitem[Knight et~al.(2022)Knight, Abd~Elbasit, and Adam]{Knight2022}
Knight, J.; Abd~Elbasit, M.A.M.; Adam, E. Land Degradation in Eritrea and
  Djibouti.
\newblock In {\em Landscapes and Landforms of the Horn of Africa: Eritrea,
  Djibouti, Somalia}; Springer International Publishing: Cham, Switzerland, 2022; pp.
  223--232.
\newblock {\url{https://doi.org/10.1007/978-3-031-05487-7_9}}.

\bibitem[Boldrocchi et~al.(2018)Boldrocchi, {Moussa Omar}, Rowat, and
  Bettinetti]{BOLDROCCHI2018812}
Boldrocchi, G.; {Moussa Omar}, Y.; Rowat, D.; Bettinetti, R.
\newblock First results on zooplankton community composition and contamination
  by some persistent organic pollutants in the Gulf of Tadjoura (Djibouti).
\newblock {\em Sci. Total. Environ.} {\bf 2018}, {\em
  627},~812--821.
\newblock {\url{https://doi.org/10.1016/j.scitotenv.2018.01.286}}.

\bibitem[Gajdzik et~al.(2021)Gajdzik, Green, Cochran, Hardenstine, Tanabe, and
  Berumen]{GAJDZIK2021112244}
Gajdzik, L.; Green, A.L.; Cochran, J.E.; Hardenstine, R.S.; Tanabe, L.K.;
  Berumen, M.L.
\newblock {Using species connectivity to achieve coordinated large-scale marine
  conservation efforts in the Red Sea}.
\newblock {\em Mar. Pollut. Bull.} {\bf 2021}, {\em 166},~112244.
\newblock {\url{https://doi.org/10.1016/j.marpolbul.2021.112244}}.

\bibitem[Billi(2022)]{Billi2022}
Billi, P. Climate Variability in the Horn of Africa Eastern Countries:
  Eritrea, Djibouti, Somalia.
\newblock In {\em Landscapes and Landforms of the Horn of Africa: Eritrea,
  Djibouti, Somalia}; Springer International Publishing: Cham, Switzerland, 2022; pp.
  1--39.
\newblock {\url{https://doi.org/10.1007/978-3-031-05487-7_1}}.

\bibitem[Seka et~al.(2024)Seka, Guo, Zhang, Han, Bayable, Ayele, Workneh,
  Bayouli, Muhirwa, and Reda]{SEKA2024141552}
Seka, A.M.; Guo, H.; Zhang, J.; Han, J.; Bayable, E.; Ayele, G.T.; Workneh,
  H.T.; Bayouli, O.T.; Muhirwa, F.; Reda, K.W.
\newblock {Evaluating the future total water storage change and hydrological
  drought under climate change over lake basins, East Africa}.
\newblock {\em J. Clean. Prod.} {\bf 2024}, {\em 447},~141552.
\newblock {\url{https://doi.org/10.1016/j.jclepro.2024.141552}}.

\bibitem[Tesfai et~al.(2021)Tesfai, Parrini, Owen-Smith, and
  Moehlman]{TESFAI2021104327}
Tesfai, R.T.; Parrini, F.; Owen-Smith, N.; Moehlman, P.D.
\newblock {African wild ass drinking behaviour on the Messir Plateau, Danakil
  Desert, Eritrea}.
\newblock {\em J. Arid. Environ.} {\bf 2021}, {\em 185},~104327.
\newblock {\url{https://doi.org/10.1016/j.jaridenv.2020.104327}}.

\bibitem[Medin et~al.(2015)Medin, Martínez-Navarro, Rivals, Libsekal, and
  Rook]{MEDIN201526}
Medin, T.; Martínez-Navarro, B.; Rivals, F.; Libsekal, Y.; Rook, L.
\newblock {The late Early Pleistocene suid remains from the
  paleoanthropological site of Buia (Eritrea): Systematics, biochronology and
  eco-geographical context}.
\newblock {\em Palaeogeogr. Palaeoclimatol. Palaeoecol.} {\bf 2015},
  {\em 431},~26--42.
\newblock {\url{https://doi.org/10.1016/j.palaeo.2015.04.020}}.

\bibitem[Ali and {Abd Ellah}(2023)]{ALI20231}
Ali, M.M.; {Abd Ellah}, R.G.
\newblock {Chapter 1-History and formation of African Lakes}. In {\em Lakes
  of Africa}; El-Sheekh, M., Elsaied, H.E., Eds.; Elsevier:  Amsterdam, The Netherlands, 
  2023; pp. 1--31.
\newblock {\url{https://doi.org/10.1016/B978-0-323-95527-0.00015-4}}.

\bibitem[Idowu and Lasisi(2020)]{IDOWU2020e00402}
Idowu, T.E.; Lasisi, K.H.
\newblock {Seawater intrusion in the coastal aquifers of East and Horn of
  Africa: A review from a regional perspective}.
\newblock {\em Sci. Afr.} {\bf 2020}, {\em 8},~e00402.
\newblock {\url{https://doi.org/10.1016/j.sciaf.2020.e00402}}.

\bibitem[Blanco-Sacristán et~al.(2022)Blanco-Sacristán, Johansen, Duarte,
  Daffonchio, Hoteit, and McCabe]{BLANCOSACRISTAN2022157098}
Blanco-Sacristán, J.; Johansen, K.; Duarte, C.M.; Daffonchio, D.; Hoteit, I.;
  McCabe, M.F.
\newblock Mangrove distribution and afforestation potential in the Red Sea.
\newblock {\em Sci. Total. Environ.} {\bf 2022}, {\em 843},~157098.
\newblock {\url{https://doi.org/10.1016/j.scitotenv.2022.157098}}.

\bibitem[Beyin(2024)]{BEYIN2024100247}
Beyin, A.
\newblock {Living by the land, gazing at the sea: Hominin occupation of
  near-coastal landscapes on the western periphery of the Red Sea}.
\newblock {\em Quat. Sci. Adv.} {\bf 2024}, {\em 16},~100247.
\newblock {\url{https://doi.org/10.1016/j.qsa.2024.100247}}.

\bibitem[Klaus(2015)]{Klaus2015}
Klaus, R. Coral Reefs and Communities of the Central and Southern Red Sea
  (Sudan, Eritrea, Djibouti, and Yemen).
\newblock In {\em The Red Sea: The Formation, Morphology, Oceanography and
  Environment of a Young Ocean Basin}; Springer: Berlin/Heidelberg,  Germany, 2015; pp. 409--451.
\newblock {\url{https://doi.org/10.1007/978-3-662-45201-1_25}}.

\bibitem[Cowburn et~al.(2019)Cowburn, Samoilys, Osuka, Klaus, Newman, Gudka,
  and Obura]{COWBURN2019182}
Cowburn, B.; Samoilys, M.; Osuka, K.; Klaus, R.; Newman, C.; Gudka, M.; Obura,
  D.
\newblock Healthy and Diverse Coral Reefs in Djibouti–A Resilient Reef
  System or Few Anthropogenic Threats?
\newblock {\em Mar. Pollut. Bull.} {\bf 2019}, {\em 148},~182--193.
\newblock {\url{https://doi.org/10.1016/j.marpolbul.2019.07.040}}.

\bibitem[Gapper et~al.(2024)Gapper, Maharjan, Li, Linstead, Tiwari, Qurban, and
  El-Askary]{GAPPER2024e38249}
Gapper, J.J.; Maharjan, S.; Li, W.; Linstead, E.; Tiwari, S.P.; Qurban, M.A.;
  El-Askary, H.
\newblock {A generalized machine learning model for long-term coral reef
  monitoring in the Red Sea}.
\newblock {\em Heliyon} {\bf 2024}, {\em 10},~e38249.
\newblock {\url{https://doi.org/10.1016/j.heliyon.2024.e38249}}.

\bibitem[Chandrasekharam et~al.(2015)Chandrasekharam, Lashin, Al~Arifi,
  Al~Bassam, and Varun]{Chandrasekharam}
Chandrasekharam, D.; Lashin, A.; Al~Arifi, N.; Al~Bassam, A.A.; Varun, C.
\newblock Evolution of geothermal systems around the Red Sea.
\newblock {\em Environ. Earth Sci.} {\bf 2015}, {\em 73},~4215--4236.
\newblock {\url{https://doi.org/10.1007/s12665-014-3710-y}}.

\bibitem[Moussa et~al.(2023)Moussa, Bayon, Dekov, Yamanaka, Shinjo, Toki, {Le
  Gall}, Grassineau, Langlade, Awaleh, and Pelleter]{MOUSSA2023104765}
Moussa, N.; Bayon, G.; Dekov, V.; Yamanaka, T.; Shinjo, R.; Toki, T.; {Le
  Gall}, B.; Grassineau, N.; Langlade, J.; Awaleh, M.;  et~al.
\newblock Mixed carbonate-siliceous hydrothermal chimneys ahead of the Asal
  propagating rift (SE Afar Rift, Republic of Djibouti).
\newblock {\em J. Afr. Earth Sci.} {\bf 2023}, {\em
  197},~104765.
\newblock {\url{https://doi.org/10.1016/j.jafrearsci.2022.104765}}.

\bibitem[Ayele et~al.(2024)Ayele, Luckett, Baptie, and Whaler]{AYELE2024105236}
Ayele, A.; Luckett, R.; Baptie, B.; Whaler, K.
\newblock The 2015 earthquake swarm in the Fentale volcanic complex (FVC): A
  geohazard risk for Ethiopia's commercial route to the Djibouti port.
\newblock {\em J. Afr. Earth Sci.} {\bf 2024}, {\em
  213},~105236.
\newblock {\url{https://doi.org/0.1016/j.jafrearsci.2024.105236}}.

\bibitem[Jalludin and Razack(1994)]{JALLUDIN1994237}
Jalludin, M.; Razack, M.
\newblock Analysis of pumping tests, with regard to tectonics, hydrothermal
  effects and weathering, for fractured Dalha and stratiform basalts, Republic
  of Djibouti.
\newblock {\em J. Hydrol.} {\bf 1994}, {\em 155},~237--250.
\newblock {\url{https://doi.org/10.1016/0022-1694(94)90167-8}}.

\bibitem[Rime et~al.(2023)Rime, Foubert, Ruch, and Kidane]{RIME2023104519}
Rime, V.; Foubert, A.; Ruch, J.; Kidane, T.
\newblock {Tectonostratigraphic evolution and significance of the Afar
  Depression}.
\newblock {\em Earth-Sci. Rev.} {\bf 2023}, {\em 244},~104519.
\newblock {\url{https://doi.org/10.1016/j.earscirev.2023.104519}}.

\bibitem[Gidafie et~al.(2024)Gidafie, Nedaw, Azagegn, Abebe, and
  Baba]{GIDAFIE2024102037}
Gidafie, D.; Nedaw, D.; Azagegn, T.; Abebe, B.; Baba, A.
\newblock {Evaluation of the source and mechanisms of groundwater recharge for
  the southern sections of the western Afar rift margin and associated rift
  floor}.
\newblock {\em J. Hydrol. Reg. Stud.} {\bf 2024}, {\em
  56},~102037.
\newblock {\url{https://doi.org/10.1016/j.ejrh.2024.102037}}.

\bibitem[Darrah et~al.(2013)Darrah, Tedesco, Tassi, Vaselli, Cuoco, and
  Poreda]{DARRAH201316}
Darrah, T.H.; Tedesco, D.; Tassi, F.; Vaselli, O.; Cuoco, E.; Poreda, R.J.
\newblock {Gas chemistry of the Dallol region of the Danakil Depression in the
  Afar region of the northern-most East African Rift}.
\newblock {\em Chem. Geol.} {\bf 2013}, {\em 339},~16--29. %MDPI: We removed ``Frontiers in Gas Geochemistry,'', please confirm % Author's reply: yes, I confirm.
\newblock 
  {\url{https://doi.org/10.1016/j.chemgeo.2012.10.036}}.

\bibitem[Varet(2022)]{Varet2022}
Varet, J. Geomorphology of Afar.
\newblock In {\em Landscapes and Landforms of the Horn of Africa: Eritrea,
  Djibouti, Somalia}; Springer International Publishing: Cham, Switzerland, 2022; pp.
  81--124.
\newblock {\url{https://doi.org/10.1007/978-3-031-05487-7_3}}.

\bibitem[Lemenkova(2022)]{Lemenkova202206}
Lemenkova, P.
\newblock {Tanzania Craton, Serengeti Plain and Eastern Rift Valley: mapping of
  geospatial data by scripting techniques}.
\newblock {\em Est. J. Earth Sci.} {\bf 2022}, {\em
  71},~61--79.
\newblock {\url{https://doi.org/10.3176/earth.2022.05}}.

\bibitem[Pinzuti et~al.(2010)Pinzuti, Mignan, and King]{PINZUTI2010169}
Pinzuti, P.; Mignan, A.; King, G.C.
\newblock Surface morphology of active normal faults in hard rock: Implications
  for the mechanics of the Asal Rift, Djibouti.
\newblock {\em Earth Planet. Sci. Lett.} {\bf 2010}, {\em
  299},~169--179.
\newblock {\url{https://doi.org/10.1016/j.epsl.2010.08.032}}.

\bibitem[Lemenkova(2022)]{Lemenkova202205}
Lemenkova, P.
\newblock {Seismicity in the Afar Depression and Great Rift Valley, Ethiopia}.
\newblock {\em Environ. Res. Eng. Manag.} {\bf 2022},
  {\em 78},~83--96.
\newblock {\url{https://doi.org/10.5755/j01.erem.78.1.29963}}.

\bibitem[Meshesha et~al.(2021)Meshesha, Chekol, and
  Negussia]{MESHESHA2021e08634}
Meshesha, D.; Chekol, T.; Negussia, S.
\newblock Major and trace element compositions of basaltic lavas from western
  margin of central main Ethiopian rift: enriched asthenosphere vs. mantle
  plume contribution.
\newblock {\em Heliyon} {\bf 2021}, {\em 7},~e08634.
\newblock {\url{https://doi.org/10.1016/j.heliyon.2021.e08634}}.

\bibitem[Mologni et~al.(2021)Mologni, Bruxelles, Schuster, Davtian, Ménard,
  Orange, Doubre, Cauliez, Taezaz, Revel, and Khalidi]{MOLOGNI2021107896}
Mologni, C.; Bruxelles, L.; Schuster, M.; Davtian, G.; Ménard, C.; Orange, F.;
  Doubre, C.; Cauliez, J.; Taezaz, H.B.; Revel, M.;  et~al.
\newblock Holocene East African monsoonal variations recorded in wave-dominated
  clastic paleo-shorelines of Lake Abhe, Central Afar region (Ethiopia \&
  Djibouti).
\newblock {\em Geomorphology} {\bf 2021}, {\em 391},~107896.
\newblock {\url{https://doi.org/10.1016/j.geomorph.2021.107896}}.

\bibitem[D'Amore et~al.(1998)D'Amore, Giusti, and Abdallah]{DAMORE1998197}
D'Amore, F.; Giusti, D.; Abdallah, A.
\newblock {Geochemistry of the High-Salinity Geothermal field of Asal, Republic
  of Djibouti, Africa}.
\newblock {\em Geothermics} {\bf 1998}, {\em 27},~197--210.
\newblock {\url{https://doi.org/10.1016/S0375-6505(97)10009-8}}.

\bibitem[Houssein and Axelsson(2010)]{HOUSSEIN2010220}
Houssein, D.E.; Axelsson, G.
\newblock Geothermal resources in the Asal Region, Republic of Djibouti: An
  update with emphasis on reservoir engineering studies.
\newblock {\em Geothermics} {\bf 2010}, {\em 39},~220--227.
\newblock {\url{https://doi.org/10.1016/j.geothermics.2010.06.006}}.

\bibitem[Awaleh and Nishijima(2024)]{AWALEH2024102894}
Awaleh, M.B.; Nishijima, J.
\newblock {Delineation of Geological Structures of Arta Geothermal Prospect in
  Djibouti Based on the Gravity Data Analysis and Interpretation}.
\newblock {\em Geothermics} {\bf 2024}, {\em 117},~102894.
\newblock {\url{https://doi.org/10.1016/j.geothermics.2023.102894}}.

\bibitem[Bouh~Houssein and Jalludin(2014)]{Houssein02112014}
Bouh~Houssein, D.; Chandrasekharam, V.C.; Jalludin, M.
\newblock {Geochemistry of Thermal Springs Around Lake Abhe, Western Djibouti}.
\newblock {\em Int. J. Sustain. Energy} {\bf 2014}, {\em
  33},~1090--1102.
\newblock {\url{https://doi.org/10.1080/14786451.2013.813027}}.

\bibitem[Castanier et~al.(1992)Castanier, Perthuisot, Rouchy, Maurin, and
  Guelorget]{CASTANIER1992811}
Castanier, S.; Perthuisot, J.P.; Rouchy, J.M.; Maurin, A.; Guelorget, O.
\newblock {Halite Ooids in Lake Asal, Djibouti: Biocrystalline Build-UPS}.
\newblock {\em Geobios} {\bf 1992}, {\em 25},~811--821.
\newblock {\url{https://doi.org/10.1016/S0016-6995(92)80063-J}}.

\bibitem[Gasse and Fontes(1989)]{GASSE198967}
Gasse, F.; Fontes, J.C.
\newblock {Palaeoenvironments and Palaeohydrology of a Tropical Closed Lake
  (Lake Asal, Djibouti) Since 10,000 yr B.P.}
\newblock {\em Palaeogeogr. Palaeoclimatol. Palaeoecol.} {\bf 1989},
  {\em 69},~67--102.
\newblock {\url{https://doi.org/10.1016/0031-0182(89)90156-9}}.

\bibitem[Dekov et~al.(2021)Dekov, Guéguen, Yamanaka, Moussa, Okumura, Bayon,
  Liebetrau, Yoshimura, Kamenov, Araoka, Makita, and Sutton]{DEKOV2021120126}
Dekov, V.M.; Guéguen, B.; Yamanaka, T.; Moussa, N.; Okumura, T.; Bayon, G.;
  Liebetrau, V.; Yoshimura, T.; Kamenov, G.; Araoka, D.;  et~al.
\newblock When a Mid-Ocean Ridge Encroaches a Continent: Seafloor-Type
  Hydrothermal Activity in Lake Asal (Afar Rift).
\newblock {\em Chem. Geol.} {\bf 2021}, {\em 568},~120126.
\newblock {\url{https://doi.org/10.1016/j.chemgeo.2021.120126}}.

\bibitem[Zan et~al.(1990)Zan, Gianelli, Passerini, Troisi, and
  Haga]{ZAN1990561}
Zan, L.; Gianelli, G.; Passerini, P.; Troisi, C.; Haga, A.O.
\newblock {Geothermal Exploration in the Republic of Djibouti: Thermal and
  Geological Data of the Hanlé and Asal Areas}.
\newblock {\em Geothermics} {\bf 1990}, {\em 19},~561--582.
\newblock {\url{https://doi.org/10.1016/0375-6505(90)90005-V}}.

\bibitem[Gong et~al.(2024)Gong, Ge, Guo, and Liu]{GONG2024149}
Gong, Z.; Ge, W.; Guo, J.; Liu, J.
\newblock {Satellite remote sensing of vegetation phenology: Progress,
  challenges, and opportunities}.
\newblock {\em ISPRS J. Photogramm. Remote Sens.} {\bf 2024},
  {\em 217},~149--164.
\newblock {\url{https://doi.org/10.1016/j.isprsjprs.2024.08.011}}.

\bibitem[Allu and Mesapam(2025)]{ALLU2025127478}
Allu, A.R.; Mesapam, S.
\newblock {Impact of remote sensing data fusion on agriculture applications: A
  review}.
\newblock {\em Eur. J. Agron.} {\bf 2025}, {\em 164},~127478.
\newblock {\url{https://doi.org/10.1016/j.eja.2024.127478}}.

\bibitem[{Chanuwan Wijesinghe} et~al.(2024){Chanuwan Wijesinghe}, {Chaminda
  Withanage}, {Kumar Mishra}, Ranagalage, Abdelrahman, and
  Fnais]{CHANUWANWIJESINGHE2024112681}
{Chanuwan Wijesinghe}, D.; {Chaminda Withanage}, N.; {Kumar Mishra}, P.;
  Ranagalage, M.; Abdelrahman, K.; Fnais, M.S.
\newblock {An application of the remote sensing derived indices for drought
  monitoring in a dry zone district, in tropical island}.
\newblock {\em Ecol. Indic.} {\bf 2024}, {\em 167},~112681.
\newblock {\url{https://doi.org/10.1016/j.ecolind.2024.112681}}.

\bibitem[Wang et~al.(2023)Wang, Sun, Cao, Wang, Zhang, and Cheng]{WANG2023311}
Wang, Y.; Sun, Y.; Cao, X.; Wang, Y.; Zhang, W.; Cheng, X.
\newblock A review of regional and Global scale Land Use/Land Cover (LULC)
  mapping products generated from satellite remote sensing.
\newblock {\em ISPRS J. Photogramm. Remote Sens.} {\bf 2023},
  {\em 206},~311--334.
\newblock {\url{https://doi.org/10.1016/j.isprsjprs.2023.11.014}}.

\bibitem[Darem et~al.(2023)Darem, Alhashmi, Almadani, Alanazi, and
  Sutantra]{DAREM2023341}
Darem, A.A.; Alhashmi, A.A.; Almadani, A.M.; Alanazi, A.K.; Sutantra, G.A.
\newblock Development of a map for land use and land cover classification of
  the Northern Border Region using remote sensing and GIS.
\newblock {\em  Egypt. J. Remote Sens. Space Sci.} {\bf
  2023}, {\em 26},~341--350.
\newblock {\url{https://doi.org/10.1016/j.ejrs.2023.04.005}}.

\bibitem[Farooq and Manocha(2024)]{FAROOQ2024101168}
Farooq, B.; Manocha, A.
\newblock Satellite-based change detection in multi-objective scenarios: A
  comprehensive review.
\newblock {\em Remote Sens. Appl. Soc. Environ.} {\bf
  2024}, {\em 34},~101168.
\newblock {\url{https://doi.org/10.1016/j.rsase.2024.101168}}.

\bibitem[Afuye et~al.(2024)Afuye, Nduku, Kalumba, Santos, Orimoloye, Ojeh,
  Thamaga, and Sibandze]{AFUYE2024103262}
Afuye, G.A.; Nduku, L.; Kalumba, A.M.; Santos, C.A.G.; Orimoloye, I.R.; Ojeh,
  V.N.; Thamaga, K.H.; Sibandze, P.
\newblock Global trend assessment of land use and land cover changes: A
  systematic approach to future research development and planning.
\newblock {\em J. King Saud Univ.-Sci.} {\bf 2024}, {\em
  36},~103262.
\newblock {\url{https://doi.org/10.1016/j.jksus.2024.103262}}.

\bibitem[Chughtai et~al.(2021)Chughtai, Abbasi, and Karas]{CHUGHTAI2021100482}
Chughtai, A.H.; Abbasi, H.; Karas, I.R.
\newblock A review on change detection method and accuracy assessment for land
  use land cover.
\newblock {\em Remote Sens. Appl. Soc. Environ.} {\bf
  2021}, {\em 22},~100482.
\newblock {\url{https://doi.org/10.1016/j.rsase.2021.100482}}.

\bibitem[Velastegui-Montoya et~al.(2023)Velastegui-Montoya, Escandón-Panchana,
  Peña-Villacreses, and Herrera-Franco]{VELASTEGUIMONTOYA2023100659}
Velastegui-Montoya, A.; Escandón-Panchana, P.; Peña-Villacreses, G.;
  Herrera-Franco, G.
\newblock Land use/land cover of petroleum activities in the framework of
  sustainable development.
\newblock {\em Clean. Eng. Technol.} {\bf 2023}, {\em
  15},~100659.
\newblock {\url{https://doi.org/10.1016/j.clet.2023.100659}}.

\bibitem[Farhan et~al.(2024)Farhan, Wu, Amin, Tariq, Guluzade, and
  Alzahrani]{FARHAN2024103689}
Farhan, M.; Wu, T.; Amin, M.; Tariq, A.; Guluzade, R.; Alzahrani, H.
\newblock Monitoring and prediction of the LULC change dynamics using time
  series remote sensing data with Google Earth Engine.
\newblock {\em Phys. Chem. Earth Parts A/B/C} {\bf 2024},
  {\em 136},~103689.
\newblock {\url{https://doi.org/10.1016/j.pce.2024.103689}}.

\bibitem[Amazirh et~al.(2023)Amazirh, Chehbouni, Bouras, Benkirane, Hssaine,
  and Entekhabi]{AMAZIRH2023104586}
Amazirh, A.; Chehbouni, A.; Bouras, E.; Benkirane, M.; Hssaine, B.; Entekhabi,
  D.
\newblock Drought cascade lag time estimation across Africa based on remote
  sensing of hydrological cycle components.
\newblock {\em Adv. Water Resour.} {\bf 2023}, {\em 182},~104586.
\newblock {\url{https://doi.org/10.1016/j.advwatres.2023.104586}}.

\bibitem[Mekonnen et~al.(2025)Mekonnen, Alamirew, Chukalla, Malede, Yalew, and
  Mengistu]{MEKONNEN20251}
Mekonnen, Y.G.; Alamirew, T.; Chukalla, A.D.; Malede, D.A.; Yalew, S.G.;
  Mengistu, A.F.
\newblock Remote sensing in hydrology: A systematic review of its applications
  in the Upper Blue Nile Basin, Ethiopia.
\newblock {\em HydroResearch} {\bf 2025}, {\em 8},~1--12.
\newblock {\url{https://doi.org/10.1016/j.hydres.2024.09.002}}.

\bibitem[Parra-Paitan et~al.(2024)Parra-Paitan, Meyfroidt, Verburg, and {zu
  Ermgassen}]{PARRAPAITAN2024103796}
Parra-Paitan, C.; Meyfroidt, P.; Verburg, P.H.; {zu Ermgassen}, E.K.
\newblock Deforestation and climate risk hotspots in the global cocoa value
  chain.
\newblock {\em Environ. Sci. Policy} {\bf 2024}, {\em 158},~103796.
\newblock {\url{https://doi.org/10.1016/j.envsci.2024.103796}}.

\bibitem[Zhao et~al.(2024)Zhao, Ciais, Wigneron, Santoro, Brandt, Kleinschroth,
  Lewis, Chave, Fensholt, Laporte, Sonwa, Saatchi, Fan, Yang, Li, Wang, Zhu,
  Xu, He, and Li]{ZHAO2024506}
Zhao, Z.; Ciais, P.; Wigneron, J.P.; Santoro, M.; Brandt, M.; Kleinschroth, F.;
  Lewis, S.L.; Chave, J.; Fensholt, R.; Laporte, N.;  et~al.
\newblock {Central African biomass carbon losses and gains during 2010–2019}.
\newblock {\em One Earth} {\bf 2024}, {\em 7},~506--519.
\newblock {\url{https://doi.org/10.1016/j.oneear.2024.01.021}}.

\bibitem[Nzabarinda et~al.(2025)Nzabarinda, Bao, Tie, Uwamahoro, Kayiranga,
  Ochege, Muhirwa, and Bao]{NZABARINDA2025114849}
Nzabarinda, V.; Bao, A.; Tie, L.; Uwamahoro, S.; Kayiranga, A.; Ochege, F.U.;
  Muhirwa, F.; Bao, J.
\newblock {Expanding forest carbon sinks to mitigate climate change in Africa}.
\newblock {\em Renew. Sustain. Energy Rev.} {\bf 2025}, {\em
  207},~114849.
\newblock {\url{https://doi.org/10.1016/j.rser.2024.114849}}.

\bibitem[Montero-Botey et~al.(2024)Montero-Botey, Kivuyo, Sitati, and
  Perea]{MONTEROBOTEY2024e03068}
Montero-Botey, M.; Kivuyo, E.; Sitati, N.; Perea, R.
\newblock {Deforestation and water availability as main drivers of
  human-elephant conflict}.
\newblock {\em Glob. Ecol. Conserv.} {\bf 2024}, {\em 54},~e03068.
\newblock {\url{https://doi.org/10.1016/j.gecco.2024.e03068}}.

\bibitem[Aksoy et~al.(2022)Aksoy, Yildirim, Gorji, Hamzehpour, Tanik, and
  Sertel]{AKSOY20221072}
Aksoy, S.; Yildirim, A.; Gorji, T.; Hamzehpour, N.; Tanik, A.; Sertel, E.
\newblock {Assessing the performance of machine learning algorithms for soil
  salinity mapping in Google Earth Engine platform using Sentinel-2A and
  Landsat-8 OLI data}.
\newblock {\em Adv. Space Res.} {\bf 2022}, {\em 69},~1072--1086.
\newblock {\url{https://doi.org/10.1016/j.asr.2021.10.024}}.

\bibitem[Nesari et~al.(2024)Nesari, Shah-Hosseini, Goodarzi, Sobhanardakani,
  and Farzaneh]{NESARI2024101989}
Nesari, H.; Shah-Hosseini, R.; Goodarzi, A.; Sobhanardakani, S.; Farzaneh, S.
\newblock {Integration of Landsat 8 (OLI) and MODIS images to monitor suspended
  particles and evaluate the spatial pattern of air pollution}.
\newblock {\em Atmos. Pollut. Res.} {\bf 2024}, {\em 15},~101989.
\newblock {\url{https://doi.org/10.1016/j.apr.2023.101989}}.

\bibitem[Kumar et~al.(2022)Kumar, Babu, Anusha, and
  Rajasekhar]{KUMAR2022100578}
Kumar, B.; Babu, K.; Anusha, B.; Rajasekhar, M.
\newblock {Geo-environmental monitoring and assessment of land degradation and
  desertification in the semi-arid regions using Landsat 8 OLI / TIRS, LST, and
  NDVI approach}.
\newblock {\em Environ. Challenges} {\bf 2022}, {\em 8},~100578.
\newblock {\url{https://doi.org/10.1016/j.envc.2022.100578}}.

\bibitem[Wright et~al.(1998)Wright, Wood, and Sylvander]{WRIGHT1998737}
Wright, D.; Wood, R.; Sylvander, B.
\newblock ArcGMT: a suite of tools for conversion between Arc/INFO® and
  Generic Mapping Tools (GMT).
\newblock {\em Comput. Geosci.} {\bf 1998}, {\em 24},~737--744.
\newblock {\url{https://doi.org/10.1016/S0098-3004(98)00067-3}}.

\bibitem[Lemenkova(2022)]{Lemenkova202217}
Lemenkova, P.
\newblock {Mapping Climate Parameters over the Territory of Botswana Using GMT
  and Gridded Surface Data from TerraClimate}.
\newblock {\em ISPRS Int. J. Geo-Inf.} {\bf 2022}, {\em
  11},~473.
\newblock {\url{https://doi.org/10.3390/ijgi11090473}}.

\bibitem[Köcher et~al.(2022)Köcher, Markaj, and Fay]{9921513}
Köcher, A.; Markaj, A.; Fay, A.
\newblock Toward a Generic Mapping Language for Transformations between RDF and
  Data Interchange Formats.
\newblock In Proceedings of the 2022 IEEE 27th International Conference on
  Emerging Technologies and Factory Automation (ETFA),  Stuttgart, Germany, 6--9 September 2022; pp. 1--4.
\newblock {\url{https://doi.org/10.1109/ETFA52439.2022.9921513}}.

\bibitem[Lemenkova(2022)]{Lemenkova202203}
Lemenkova, P.
\newblock {Console-Based Mapping of Mongolia Using GMT Cartographic Scripting
  Toolset for Processing TerraClimate Data}.
\newblock {\em Geosciences} {\bf 2022}, {\em 12},~1--36.
\newblock {\url{https://doi.org/10.3390/geosciences12030140}}.

\bibitem[Bawa et~al.(2024)Bawa, Onotu, Akomolafe, and Sa’i]{10537274}
Bawa, S.; Onotu, A.A.; Akomolafe, E.A.; Sa’i, U.I.
\newblock Understanding the Topography of the Gulf of Guinea Seabed Using GMT
  Scripting and GEBCO Gridded Data.
\newblock In Proceedings of the 2024 IEEE Mediterranean and Middle-East
  Geoscience and Remote Sensing Symposium (M2GARSS),  Oran, Algeria, 15--17 April 2024; pp. 516--520.
\newblock {\url{https://doi.org/10.1109/M2GARSS57310.2024.10537274}}.

\bibitem[Tangkitjaroenmongkol et~al.(2011)Tangkitjaroenmongkol, Kaittisin, and
  Ongwattanakul]{5930121}
Tangkitjaroenmongkol, R.; Kaittisin, S.; Ongwattanakul, S.
\newblock Inbound logistics cassava starch planning: With application of GIS
  and K-means clustering methods in Thailand.
\newblock In Proceedings of the 2011 Eighth International Joint Conference on
  Computer Science and Software Engineering (JCSSE),  Nakhonpathom, Thailand, 11--13 May 2011; pp. 204--209.
\newblock {\url{https://doi.org/10.1109/JCSSE.2011.5930121}}.

\bibitem[Xing et~al.(2023)Xing, Zhao, Song, Liu, Chai, Zhang, Li, Jia, Zhao,
  and Xu]{10430357}
Xing, F.; Zhao, Y.; Song, Z.; Liu, Z.; Chai, W.; Zhang, Y.; Li, M.; Jia, T.;
  Zhao, J.; Xu, M.
\newblock GIS Gas Analysis Training Program Based on K-Means Clustering Method
  Training Auxiliary Analysis.
\newblock In Proceedings of the 2023 International Conference on Educational
  Knowledge and Informatization (EKI), Guangzhou, China, 22--24 September  2023; pp. 78--82.
\newblock {\url{https://doi.org/10.1109/EKI61071.2023.00025}}.

\bibitem[Lemenkova(2023)]{Lemenkova202315}
Lemenkova, P.
\newblock {Mapping Wetlands of Kenya Using Geographic Resources Analysis
  Support System (GRASS GIS) with Remote Sensing Data}.
\newblock {\em Transylv. Rev. Syst. Ecol. Res.}
  {\bf 2023}, {\em 25},~1--18.
\newblock {\url{https://doi.org/10.2478/trser-2023-0008}}.

\bibitem[Huo et~al.(2011)Huo, Moon, Lee, Seung, and Kwon]{5739741}
Huo, X.J.; Moon, K.S.; Lee, S.H.; Seung, T.Y.; Kwon, K.R.
\newblock Protecting GIS vector map using the k-means clustering algorithm and
  odd-even coding.
\newblock In Proceedings of the 2011 17th Korea-Japan Joint Workshop on
  Frontiers of Computer Vision (FCV), Ulsan, Republic of Korea, 9--11 February 2011; pp. 1--5.
\newblock {\url{https://doi.org/10.1109/FCV.2011.5739741}}.

\bibitem[Chowdhury(2024)]{CHOWDHURY2024100800}
Chowdhury, M.S.
\newblock Comparison of accuracy and reliability of random forest, support
  vector machine, artificial neural network and maximum likelihood method in
  land use/cover classification of urban setting.
\newblock {\em Environ. Challenges} {\bf 2024}, {\em 14},~100800.
\newblock {\url{https://doi.org/10.1016/j.envc.2023.100800}}.

\bibitem[Lemenkova(2023)]{Lemenkova202323}
Lemenkova, P.
\newblock {Using Open-Source Software GRASS GIS for Analysis of the
  Environmental Patterns in Lake Chad, Central Africa}.
\newblock {\em Die Bodenkultur J. Land Manag. Food Environ.} {\bf 2023}, {\em 74},~49--64.
\newblock {\url{https://doi.org/10.2478/boku-2023-0005}}.

\bibitem[Pedregosa et~al.(2018)Pedregosa, Varoquaux, Gramfort, Michel, Thirion,
  Grisel, Blondel, Müller, Nothman, Louppe, Prettenhofer, Weiss, Dubourg,
  Vanderplas, Passos, Cournapeau, Brucher, Perrot, and Édouard
  Duchesnay]{pedregosa2018}
Pedregosa, F.; Varoquaux, G.; Gramfort, A.; Michel, V.; Thirion, B.; Grisel,
  O.; Blondel, M.; Müller, A.; Nothman, J.; Louppe, G.;  et~al.
\newblock Scikit-learn: Machine Learning in Python. \emph{ arXiv}  \textbf{2018},
 % \href{http://arxiv.org/abs/1201.0490}{{\normalfont [arXiv:cs.LG/1201.0490]}}.
\newblock {\url{https://doi.org/10.48550/arXiv.1201.0490}}.

\bibitem[Wenting et~al.(2004)Wenting, Bingfang, Yichen, and Yuan]{1370179}
Wenting, X.; Bingfang, W.; Yichen, T.; Yuan, Z.
\newblock Mapping plant diversity of broad-leaved forest ecosystem using
  Landsat TM data.
\newblock In Proceedings of the IGARSS 2004. 2004 IEEE International Geoscience
  and Remote Sensing Symposium,  Anchorage, AK,  USA, 20--24 September 2004; Volume~7, pp. 4598--4600.
\newblock {\url{https://doi.org/10.1109/IGARSS.2004.1370179}}.

\bibitem[Dobrowski et~al.(2008)Dobrowski, Safford, Cheng, and Ustin]{Dobrowski}
Dobrowski, S.Z.; Safford, H.D.; Cheng, Y.B.; Ustin, S.L.
\newblock {Mapping Mountain Vegetation using Species Distribution Modeling,
  Image-Based Texture Analysis, and Object-Based Classification}.
\newblock {\em Appl. Veg. Sci.} {\bf 2008}, {\em 11},~499--508.
\newblock {\url{https://doi.org/10.3170/2008-7-18560}}.

\bibitem[Moore et~al.(2007)Moore, Hoffmann, and Glenn]{Moore}
Moore, C.A.; Hoffmann, G.A.; Glenn, N.F.
\newblock {Quantifying Basalt Rock Outcrops in NRCS Soil Map Units Using
  Landsat-5 Data}.
\newblock {\em Soil Surv. Horizons} {\bf 2007}, {\em 48},~59--62.
\newblock {\url{https://doi.org/10.2136/sh2007.3.0059}}.

\bibitem[Masemola et~al.(2019)Masemola, Cho, and Ramoelo]{8688643}
Masemola, C.; Cho, M.A.; Ramoelo, A.
\newblock Assessing the Effect of Seasonality on Leaf and Canopy Spectra for
  the Discrimination of an Alien Tree Species, Acacia Mearnsii, From
  Co-Occurring Native Species Using Parametric and Nonparametric Classifiers.
\newblock {\em IEEE Trans. Geosci. Remote Sens.} {\bf 2019},
  {\em 57},~5853--5867.
\newblock {\url{https://doi.org/10.1109/TGRS.2019.2902774}}.

\bibitem[Lemenkova(2023{\natexlab{a}})]{Lemenkova202322}
Lemenkova, P.
\newblock {Monitoring Seasonal Fluctuations in Saline Lakes of Tunisia Using
  Earth Observation Data Processed by GRASS GIS}.
\newblock {\em Land} {\bf 2023}, {\em 12},~1995.
\newblock {\url{https://doi.org/10.3390/land12111995}}.

\bibitem[Lemenkova(2023{\natexlab{b}})]{Lemenkova202313}
Lemenkova, P.
\newblock {A GRASS GIS Scripting Framework for Monitoring Changes in the
  Ephemeral Salt Lakes of Chotts Melrhir and Merouane, Algeria}.
\newblock {\em Appl. Syst. Innov.} {\bf 2023}, {\em 6},~61.
\newblock {\url{https://doi.org/10.3390/asi6040061}}.

\bibitem[Eldosouky et~al.(2021)Eldosouky, Pham, El-Qassas, Hamimi, and
  Oksum]{Eldosouky2021}
Eldosouky, A.M.; Pham, L.T.; El-Qassas, R.A.Y.; Hamimi, Z.; Oksum, E.
  Lithospheric Structure of the Arabian--Nubian Shield Using Satellite
  Potential Field Data.
\newblock In {\em The Geology of the Arabian-Nubian Shield}; Springer
  International Publishing: Cham, Switzerland, 2021; pp. 139--151.
\newblock {\url{https://doi.org/10.1007/978-3-030-72995-0_6}}.

\end{thebibliography}
\PublishersNote{}

%%%%%%%%%%%%%%%%%%%%%%%%%%%%%%%%%%%%%%%%%%
\end{adjustwidth}
\end{document}
